% ulm-lam.tex -- design sketch: solving the Linear Algebra Mapping Problem
% with MLIR.  Build: latexmk (LuaLaTeX, see latexmkrc).
\documentclass[11pt,a4paper]{article}

\usepackage{fontspec}
\usepackage{unicode-math}
\setmainfont{Libertinus Serif}
\setsansfont{Libertinus Sans}
\setmathfont{Libertinus Math}
\setmonofont{JuliaMono}[Scale=MatchLowercase]

\usepackage[margin=2.5cm]{geometry}
\usepackage{microtype}
\usepackage{booktabs,tabularx,array,multirow,graphicx}
\usepackage[inline]{enumitem}
\usepackage[skip=\medskipamount]{parskip}   % no indent, \medskip between paragraphs
\usepackage{xcolor,luacolor,lua-ul}
\usepackage{listings}
\usepackage{caption}   % \captionof for listings that break across pages
\usepackage{tikz}
\usetikzlibrary{positioning,arrows.meta,fit,backgrounds}
\usepackage[hidelinks]{hyperref}
\usepackage[nameinlink,capitalise]{cleveref}

\setlist{itemsep=2pt,topsep=4pt}
\renewcommand{\arraystretch}{1.15}
\setlength{\emergencystretch}{3em}

% ---------------------------------------------------------------- notation
\newcommand{\T}{^{\mkern-1.5mu\mathsf{T}}}
\newcommand{\inv}{^{-1}}
\newcommand{\invT}{^{-\mathsf{T}}}

% ---------------------------------------------------------------- listings
\definecolor{kw}{HTML}{1F4E9A}
\definecolor{cm}{HTML}{6B6B6B}
\definecolor{st}{HTML}{8A3B12}
\definecolor{ty}{HTML}{2E6B30}
\definecolor{bg}{gray}{0.95}      % code blocks
\definecolor{ttbg}{gray}{0.92}    % inline typewriter text
% Grey background behind \texttt via lua-ul leaders: no effect on line
% spacing, line breaks still possible.
\let\origtexttt\texttt
\renewcommand{\texttt}[1]{\highLight[ttbg]{\origtexttt{#1}}}

\lstdefinelanguage{mlir}{
  alsoletter={.,_,\%,\#,!,@},
  morekeywords={func.func,return,func.call},
  morekeywords=[2]{ulm_lam.mul,ulm_lam.add,ulm_lam.scale,ulm_lam.transpose,
    ulm_lam.inverse,ulm_lam.identity,ulm_lam.assume,ulm_lamk.gemm,ulm_lamk.symm,
    ulm_lamk.trmm,ulm_lamk.trsm,ulm_lamk.syrk,ulm_lamk.gemv,ulm_lamk.trsv,ulm_lamk.potrf,
    ulm_lamk.getrf,ulm_lamk.getrs},
  morekeywords=[3]{tensor,memref,f64,i32,f32,index,strided},
  morecomment=[l]{//},
  morestring=[b]",
  sensitive=true}
\lstdefinelanguage{tablegen}{
  morekeywords={def,class,let,include,defvar,foreach,in,if,then,else,
    dag,string,list,code,bit,int,bits,ifndef,define,endif},
  morekeywords=[2]{ins,outs},
  morecomment=[l]{//},
  morestring=[b]",
  sensitive=true}
\lstdefinestyle{base}{
  basicstyle=\ttfamily\footnotesize,
  keywordstyle=\color{kw}\bfseries,
  keywordstyle=[2]\color{kw},
  keywordstyle=[3]\color{ty},
  commentstyle=\color{cm}\itshape,
  stringstyle=\color{st},
  backgroundcolor=\color{bg},
  frame=none,
  xleftmargin=4pt, xrightmargin=4pt,
  framexleftmargin=4pt, framexrightmargin=4pt,
  aboveskip=6pt, belowskip=6pt,
  columns=fullflexible, keepspaces=true,
  showstringspaces=false, breaklines=true,
  upquote=true, tabsize=8}
\lstdefinestyle{mlir}{style=base,language=mlir}
\lstdefinestyle{td}{style=base,language=tablegen}
\lstdefinestyle{sh}{style=base,language=bash,morekeywords={}}
\lstdefinestyle{c}{style=base,language=C}
\lstdefinestyle{py}{style=base,language=Python,morekeywords={with,as}}
\lstset{style=mlir}

% Excerpts from the real sources (see ulm-lam-doc.lua).
\directlua{ulmlamdoc = dofile("ulm-lam-doc.lua")}
\newcommand{\code}[1]{\texttt{\directlua{ulmlamdoc.code("\luaescapestring{\detokenize{#1}}")}}}
\let\file\code
\newcommand{\srcdef}[3][]{%
  \directlua{ulmlamdoc.def("\luaescapestring{#2}", "\luaescapestring{#3}",
                         "\jobname-#3.lst")}%
  \lstinputlisting[style=td,#1]{\jobname-#3.lst}}
\newcommand{\srcbetween}[5][]{%
  \directlua{ulmlamdoc.between("\luaescapestring{#2}", "\luaescapestring{#3}",
                             "\luaescapestring{#4}", "\jobname-#5.lst")}%
  \lstinputlisting[#1]{\jobname-#5.lst}}

\newcommand{\src}{../}

% ---------------------------------------------------------------- document
\title{ulm.lam: Write the Formula,\\Get the Expert's BLAS and LAPACK Code\\[4pt]
  \large Ulm Linear Algebra Mapping with MLIR and ulmBLAS}
\author{Michael Lehn}
\date{\today}

\begin{document}

\maketitle

\begin{abstract}
\noindent
You write linear algebra as on paper, in Python or Julia, with inverses,
transposes and matrix properties; ulm.lam compiles it into the code an
expert in numerical linear algebra would write on top of BLAS (Basic
Linear Algebra Subprograms) and LAPACK (Linear Algebra PACKage):
factorisations instead of inverses, computed once and pulled out of
loops, products parenthesised by cost, symmetric results computed as one
triangle, no memory allocated inside nested calls, and errors reported as
LAPACK's \code{info}.
A training loop of a neural-network layer, written as formulas, runs as
fast as PyTorch on the same BLAS, and $5\times$ faster than PyTorch's
autograd for a small layer.
The prediction of a Gaussian process, five lines of Python, runs
$1.4\times$ faster than the same steps written by hand with SciPy and
$8\times$ faster than the formula written in NumPy, and scales to
$n = 100\,000$ on a cluster node, where it is as fast as the same steps
written by hand with Intel's MKL (Math Kernel Library).
A loop of Kalman measurement updates, written as the formula, is
$3\times$ faster than NumPy code in which the operator was formed by hand
before the loop.
In Julia, a \code{@lam} function is called with Julia arrays, runs on
Julia's own BLAS, and \code{@code_lam} shows the BLAS and LAPACK calls
ulm.lam derived.
Behind it is a solution of the \emph{Linear Algebra Mapping Problem}
(LAMP) on MLIR (Multi-Level Intermediate Representation)~\cite{mlir}:
two dialects (the algebra with matrix properties, and the kernels in
destination-passing style), kernel selection by a generalised matrix
chain algorithm in the spirit of Linnea~\cite{linnea-toms,linnea-thesis},
MLIR's bufferization for the memory, and a lowering to C++ calling
ulmBLAS~\cite{ulmblas}, or any CBLAS and LAPACK.
\end{abstract}

\setcounter{tocdepth}{1}
\tableofcontents

% =====================================================================
\section{ulm.lam at a glance}\label{sec:glance}
% =====================================================================

\paragraph{What you write.}
A neural-network layer trained by gradient descent: a linear layer
(PyTorch's \code{nn.Linear}, $Y \approx XW + \mathbf{1}b\T$), the forward
pass, the gradient of the mean squared error and the update step written
as formulas, the whole training loop one ulm.lam function:
\begin{lstlisting}[style=py]
from ulm import lam
from ulm.lam import Float, Index, Matrix, Vector

batch_x = lam.extern("batch_x", [Index], Matrix("B", "p"))   # Python callbacks
batch_y = lam.extern("batch_y", [Index], Matrix("B", "k"))

@lam.function
def train(W0: Matrix("p", "k"), b0: Vector("k"), ones: Vector("B"),
          step: Float, steps: Index):
    with lam.loop(steps, carry=[W0, b0]) as L:
        W, b = L.carried
        X, Y = batch_x(L.index), batch_y(L.index)
        R = X @ W + ones @ b.T - Y                      # forward, residual
        L.next(W - step * (X.T @ R), b - step * (R.T @ ones))
    return L.results
\end{lstlisting}
ulm.lam turns each step into four BLAS calls and allocates nothing in the
loop: $XW - Y$ by one \code{gemm} into a copy of $Y$ ($\beta = -1$), the
bias as a rank-1 \code{gemm}, the update $W - \mathrm{step}\,X\T R$ by one
\code{gemm} with $\alpha = -\mathrm{step}$, $\beta = 1$ directly into $W$,
and $b$ by one \code{gemv}.
On the same BLAS it is as fast as PyTorch, and $5\times$ faster than
PyTorch's autograd for a small layer (\cref{sec:torch}).

\paragraph{A formula.}
The prediction of a Gaussian process (GP), the textbook formula
\[
  \mu = K_*\T (K + \sigma^2 I)\inv y, \qquad
  \Sigma = K_{**} - K_*\T (K + \sigma^2 I)\inv K_*,
\]
written in Python:
\begin{lstlisting}[style=py]
from ulm import lam
from ulm.lam import Float, Matrix, Vector

@lam.function
def gp(K: Matrix("n", "n", "symmetric"), s2: Float, Ks: Matrix("n", "m"),
       Kss: Matrix("m", "m", "symmetric"), y: Vector("n")):
    Kn = lam.assume(K + s2 * lam.eye("n"), "spd")
    return Ks.T @ Kn.inv @ y, Kss - Ks.T @ Kn.inv @ Ks

lib = lam.compile(gp)                  # C++ calling ulmBLAS, compiled once
mean, cov = lib.gp(K, 0.01, Ks, Kss, y)  # NumPy arrays in and out
\end{lstlisting}
The sizes are names: one compiled library serves all $n$ and $m$.
\code{lam.assume} states what the compiler cannot know ($K + \sigma^2 I$
is positive definite for $\sigma^2 > 0$); the Cholesky factorisation
checks it at run time.

\paragraph{How fast it is.}
Bike-sharing data (UCI, University of California, Irvine), $m = 1000$
test points, Apple M5 (seconds, without computing the kernel matrices):

\begin{center}
\begin{tabular}{rrrrrr}
\toprule
$n$ & ulm.lam/ulmBLAS & ulm.lam/OpenBLAS & SciPy by hand & scikit-learn & NumPy literal \\
\midrule
2000 & 0.036 & 0.046 & 0.053 & 0.119 & 0.194 \\
4000 & 0.133 & 0.189 & 0.242 & 0.547 & 1.109 \\
8000 & 0.886 & 1.105 & 1.209 & 2.970 & 6.983 \\
\bottomrule
\end{tabular}
\end{center}

\noindent
``SciPy by hand'' is what scikit-learn does inside, written directly with
SciPy; ``NumPy literal'' is the formula with \code{np.linalg.inv};
scikit-learn's time includes the kernel matrices, which it computes
inside (0.78\,s of its 2.97\,s at $n = 8000$).
On the same BLAS (OpenBLAS, the one NumPy uses here) the generated code is
faster than the hand-written version; on ulmBLAS it is faster still.
On a cluster node (2$\times$ AMD EPYC 9454, 96 cores) the same function
predicts terrain altitudes from $n = 100\,000$ road points
(\cref{sec:gpcluster}): the factorisation of an 80\,GB matrix, in
154\,s on MKL, as fast as the same steps written by hand with MKL.
And in a loop, the Kalman measurement update
$b_i = (H\T R\inv H)\inv H\T R\inv y_i$ for 200 measurements
($n = 4000$) takes 122\,ms written as the formula, against 368\,ms for
NumPy code in which the operator was formed by hand before the loop, and
270\,s for the formula written in NumPy (\cref{sec:python}).

\paragraph{The same in Julia.}
The Julia package (\cref{sec:julia}) uses Julia's notation; Julia's
column-major arrays are passed without copies:
\begin{lstlisting}[style=py]
using UlmLam, LinearAlgebra

@lam function gp(K::Mat(:n, :n, :symmetric), s2::Float, Ks::Mat(:n, :m),
                 Kss::Mat(:m, :m, :symmetric), y::Vec(:n))
    Kn = assume(K + s2 * I, :spd)
    (Ks' * inv(Kn) * y, Kss - Ks' * inv(Kn) * Ks)
end

lib = UlmLam.compile(gp)
mean, cov = lib.gp(K, 0.01, Ks, Kss, y)
\end{lstlisting}
Compared with what a Julia programmer writes by hand
(\code{C = cholesky(Symmetric(K + s2*I))}, \code{C \textbackslash{} y},
\code{V = C.L \textbackslash{} Ks}, \code{Kss - V'V}) and with the formula
written literally (\code{inv}), RBF (radial basis function) kernel
matrices of random points, $m = 1000$, Apple M5, Julia 1.13 with its
OpenBLAS on 4 threads (its default here; seconds):

\begin{center}
\begin{tabular}{rrrrr}
\toprule
$n$ & ulm.lam/ulmBLAS & ulm.lam/OpenBLAS & Julia by hand & Julia literal \\
\midrule
2000 & 0.037 & 0.045 & 0.071 & 0.212 \\
4000 & 0.137 & 0.216 & 0.332 & 1.359 \\
8000 & 0.907 & 1.462 & 1.922 & 8.828 \\
\bottomrule
\end{tabular}
\end{center}

\noindent
In the Kalman loop (200 steps, $n = 4000$), ulm.lam takes about 140\,ms,
Julia with the operator formed by hand before the loop 390 to 550\,ms
(it varies between runs), and the formula written in Julia 400 to 520\,s
(\file{julia/examples/kalman.jl}).

\paragraph{Why it is fast.}
ulm.lam is a compiler: the formula is traced into an intermediate
representation, mapped onto BLAS and LAPACK calls, and compiled to C++.
The mapping makes the decisions an expert makes when writing such code by
hand:
\begin{itemize}
  \item inverses become factorisations and triangular solves; a matrix is
        factorised once, however often its inverse appears;
  \item properties are exploited: Cholesky instead of LU for positive
        definite matrices, \code{syrk} for symmetric products (half the
        flops of \code{gemm}), one triangle of symmetric matrices;
  \item chains of products are parenthesised to minimise the cost,
        at run time if the best order depends on the sizes;
  \item factorisations and other loop-invariant work are pulled out of
        loops, and the whole operator is formed before a loop if the
        number of steps makes it pay off;
  \item no memory is allocated dynamically, not even inside nested
        calls: all temporaries live in one workspace with a LAPACK-style
        size query, kernels work in place where an input is no longer
        needed, and errors are returned as LAPACK's \code{info}.
\end{itemize}
In short: it generates the code a programmer experienced in numerical
linear algebra would write on top of BLAS and LAPACK.

\paragraph{What the mapping does.}
With $K + \sigma^2 I = L L\T$ (Cholesky, $L$ lower triangular),
$(K + \sigma^2 I)\inv = L\invT L\inv$, so
\[
  \mu = K_*\T \alpha \quad\text{with}\quad \alpha = L\invT (L\inv y),
  \qquad
  \Sigma = K_{**} - V\T V \quad\text{with}\quad V = L\inv K_*.
\]
One factorisation serves both results, no inverse is formed, and $V\T V$
is symmetric, so only one triangle is computed.
As calls (all temporaries in the workspace $w$):
\begin{center}
\small
\begin{tabular}{@{}lll@{}}
\toprule
& operation & kernel \\
\midrule
1 & $W \leftarrow K$, \ $W \leftarrow W + \sigma^2 I$ & copy, \code{adddiag} \\
2 & $W = L L\T$ ($L$ overwrites $W$); if it fails: return \code{info} 1 & \code{potrf} \\
3 & $a \leftarrow y$, \ $a \leftarrow L\inv a$, \ $a \leftarrow L\invT a$ & copy, \code{trsv}, \code{trsv} \\
4 & $\mu \leftarrow K_*\T a$ & \code{gemv} \\
5 & $V \leftarrow K_*$, \ $V \leftarrow L\inv V$ & copy, \code{trsm} \\
6 & $\Sigma \leftarrow K_{**}$, \ $\Sigma \leftarrow \Sigma - V\T V$ (lower triangle), mirror it & copy, \code{syrk}, complete \\
\bottomrule
\end{tabular}
\end{center}
\noindent
The copies exist because $K$, $K_*$ and $K_{**}$ are \code{in}
parameters, which the function must not overwrite; declared
\code{InOut}, $K$ is factorised in place (\cref{sec:gpcluster}).
\Cref{lst:gpcpp} is the C++ that ulm.lam generates from the five lines of
Python.

% (not a float: the listing is longer than a page and breaks)
\begin{flushleft}
\IfFileExists{\src python/examples/generated/gp.cpp}{%
  \srcbetween[style=c]{\src python/examples/generated/gp.cpp}{/* gp: matrices}{}{gpcpp}}{%
  \fbox{run \code{make gp} to generate this listing}}
\captionsetup{hypcap=false}
\captionof{figure}{The function \code{gp} generated from the Python
above (\file{python/examples/generated/gp.cpp}, without the
\code{\#include} lines and the helpers \code{ulm_lam_copy},
\code{ulm_lam_adddiag} and \code{ulm_lam_sycomplete} in front of it).
The interface comment says what the caller passes and what \code{info}
means; the outer \code{if}s implement the workspace query, the inner ones
\code{info}; then come the steps of the table above.}
\label{lst:gpcpp}
\end{flushleft}

% =====================================================================
\section{The linear algebra mapping problem}
% =====================================================================

A linear algebra problem is written as mathematics:
\begin{equation}\label{eq:gls}
  b \;=\; (X\T M\inv X)\inv X\T M\inv y,
  \qquad M \text{ symmetric positive definite (SPD)},\; X \text{ of full rank},
\end{equation}
the generalised least-squares estimator.
It is executed as a sequence of calls to tuned kernels:
\code{potrf}, \code{trsm}, \code{syrk}, \code{gemv}, \code{trsv}.
Choosing that sequence is the LAMP~\cite{lamp-paper}.
The choice includes parenthesisations, which inverses become solves,
which factorisations to compute, which properties to exploit, and which
subexpressions to share.
The problem is NP-complete (NP: nondeterministic polynomial time) in
general~\cite{lamp-paper}.
Mainstream languages and libraries (Matlab, Julia, Eigen, Armadillo, NumPy,
TensorFlow, PyTorch) solve it only partially~\cite{lamp-paper,tfpt}.

Linnea~\cite{linnea-toms,linnea-thesis} is the reference solution.
It normalises the input by term rewriting with MatchPy~\cite{matchpy}, infers
matrix properties, and runs a best-first search over a graph of
derivations.
It emits Julia code that calls BLAS/LAPACK and reaches expert-level
performance.
It is also a Python research prototype, last changed in 2021.

\paragraph{Why MLIR.}
MLIR provides what Linnea had to build by hand, and much that it lacks:
\begin{itemize}
  \item a typed, verified intermediate representation (IR) and a textual format that are stable;
  \item rewrite infrastructure, common subexpression elimination (CSE), and \code{lit}/\code{FileCheck}
        testing;
  \item One-Shot Bufferize, which decides in-place updates and copies.
        This replaces Linnea's hand-written ``overwriting'' logic;
  \item lowering to LLVM (originally ``Low Level Virtual Machine'', today
        just the project's name), and a code-generation fallback (\code{linalg})
        when no library kernel fits;
  \item frontends, Python bindings and an ecosystem that is maintained.
\end{itemize}

\paragraph{What MLIR does not provide.}
The core of the LAMP is not covered by any upstream component:
\begin{itemize}
  \item \textbf{search instead of greedy rewriting.}
        MLIR's pattern driver applies rewrites greedily.
        Kernel selection must compare alternatives under a cost model;
  \item \textbf{associative-commutative matching.}
        MLIR patterns match DAGs (directed acyclic graphs) structurally, while kernels have to be
        found inside products and sums of any length;
  \item \textbf{a dialect at the right level.}
        \code{linalg} has no inverse, no notion of SPD, and no transpose
        as an algebraic operator.
\end{itemize}
The design below keeps these parts explicit and small, and uses upstream
MLIR for everything else.

\paragraph{What this repository contains.}
\Cref{tab:files} lists the files.
The ODS files are \emph{checked}: \file{scripts/check-ods.sh} runs every
\code{mlir-tblgen} backend the dialect needs, including validation of the
custom assembly formats.
The dialects are \emph{implemented}: \file{lib/} contains both dialects,
generated from the ODS, and a first set of passes, built as a plugin for
\code{mlir-opt} (\cref{sec:impl}).
The worked example is partly automatic.
Its stages~1–4 are written by hand as a design: they show Linnea's
derivation.
They parse and verify with the implemented dialects, and stage~3
bufferizes into exactly the copies that stage~4 shows.
From stage~2 on, the passes produce C++ calling ulmBLAS by themselves;
it is compiled, run and verified numerically (\cref{sec:example}).
The loop example (\cref{sec:loops}) and eight kernel selection cases with
inverses, solves and sums are compiled the same way and checked against
naive evaluations.

\begin{table}[t]
\centering\small
\begin{tabularx}{\linewidth}{@{}lX@{}}
\toprule
File & Content \\
\midrule
\file{include/ulm_lam/UlmLamBase.td} & the \code{ulm_lam} and \code{ulm_lamk} dialects \\
\file{include/ulm_lam/UlmLamAttrs.td} & property encoding \code{#ulm_lam.props}, kernel flags \\
\file{include/ulm_lam/UlmLamInterfaces.td} & property inference and cost interfaces \\
\file{include/ulm_lam/UlmLamOps.td} & \code{ulm_lam} operations \\
\file{include/ulm_lam/UlmLamKernelOps.td} & \code{ulm_lamk} operations (BLAS/LAPACK) \\
\file{include/ulm_lam/UlmLamPasses.td} & passes of the pipeline \\
\file{include/ulm_lam/UlmLamKernels.td} & declarative kernel catalogue \\
\file{examples/gls/0[1-4]-*.mlir} & the worked example: design stages 1–4 \\
\file{examples/gls/run.sh}, \file{driver.cc} & from stage 2 to C++: compile, run, verify \\
\file{examples/gls/gls_inout.mlir} & GLS with $M$, $y$ declared \code{inout} \\
\file{examples/cases/} & inverses, solves and sums, checked numerically \\
\file{examples/kalman/} & formula (1) in a loop: changing $y$; covariance updated in place (\cref{sec:kalman}) \\
\file{examples/compose/} & workspaces of functions that call functions with workspaces \\
\file{examples/*/generated/} & the C++ generated from each example (\cref{app:generated}) \\
\file{examples/loop/} & temporaries in loops (\cref{sec:loops}) \\
\file{lib/UlmLamResultsToOutParams.cpp} & implemented pass: matrices by reference \\
\file{lib/ConvertUlmLamkToUlmBLAS.cpp} & lowering to C++ calling ulmBLAS \\
\file{lib/UlmLamWorkspace.cpp} & temporaries from a caller-provided workspace \\
\file{examples/workspace.h} & the workspace convention for the C++ drivers \\
\file{python/ulm/lam/} & the Python frontend and runtime (\cref{sec:python}) \\
\file{python/examples/}, \file{python/tests/} & Python example and unit tests \\
\file{scripts/ulm-lam-compile.sh} & the pipeline: \code{ulm_lam} ops to C++ (all stages) \\
\file{scripts/build-plugin.sh} & build the \code{mlir-opt} pass plugin \\
\file{scripts/check-ods.sh} & run all \code{mlir-tblgen} backends \\
\file{scripts/check-stages.sh} & stages 1–4 parse and verify; stage 3 bufferizes as stage 4 \\
\file{lib/UlmLamDialect.{h,cpp}}, \file{UlmLamkDialect.cpp} & the dialects: inference, verifiers, patterns, costs, bufferization \\
\file{lib/UlmLamPasses.cpp} & normal form, property inference, kernel selection, property stripping \\
\file{lib/Plugin.cpp} & \code{mlir-opt} dialect and pass plugin entry points \\
\file{doc/} & this document (LuaLaTeX) \\
\bottomrule
\end{tabularx}
\caption{Repository contents. Toolchain used: Homebrew LLVM/MLIR 23.1.4,
ulmBLAS (commit c2273e7, November 2026).}
\label{tab:files}
\end{table}

% =====================================================================
\section{Related work}
% =====================================================================

\paragraph{Linnea.}
Linnea~\cite{linnea-toms,linnea-thesis} works in four steps:
\begin{enumerate}
  \item it rewrites the input into a normal form, with associative
        operators flattened, transposes and inverses pushed to the
        leaves, and properties propagated;
  \item it represents the program as a graph of derivations, where each
        edge applies a kernel or a factorisation;
  \item it explores that graph best-first, with FLOPs (floating-point
        operations) as cost.
        The search is \emph{anytime}: a first solution typically comes in
        under a second, and better ones follow if it runs longer;
  \item it finds matrix products with the generalised matrix chain (GMC)
        algorithm~\cite{gmc}.
\end{enumerate}
GMC is a dynamic programme over parenthesisations.
It chooses a kernel for every binary product and takes transposes,
inverses and properties into account.
Matching uses MatchPy's associative-commutative matching with sequence
variables~\cite{matchpy}.
The implementation is Python.
Its kernel definitions exist for Julia and C, and it generates Julia
(see \url{https://github.com/HPAC/linnea}).

\paragraph{MOM / MLIR-Linnea.}
Chelini et al.~\cite{mom} proposed a matrix dialect for MLIR that carries
properties in the types, has n-ary multiplication, and contains a
generalised matrix-chain pass.
It lowers to \code{linalg}.
The prototype\footnote{\url{https://github.com/Algebraic-Programming/MLIR-Linnea},
last commit June 2022.} has property propagation and matrix-chain
reordering.
Kernel detection is stubbed: \code{isTRSM()} returns \code{false}.
It has no search, and it is pinned to a 2022 LLVM.
This sketch follows MOM in putting properties into types.
It differs in three ways:
\begin{itemize}
  \item it uses the builtin tensor type with an encoding, as
        \code{sparse_tensor} does, instead of a new type;
  \item it adds a kernel dialect in destination-passing style, so that
        upstream bufferization can do the memory planning;
  \item it makes the search a separate, explicit component.
\end{itemize}

\paragraph{Equality saturation.}
Equality saturation (egg~\cite{egg}, egglog~\cite{egglog}) explores all
rewrites at once and extracts a cheapest term.
SPORES~\cite{spores} applies it to linear algebra in a relational setting.
Two MLIR integrations exist:
\begin{itemize}
  \item the \code{eqsat} dialect~\cite{eqsat}, which represents an
        e-graph inside MLIR and is driven by PDL (Pattern Description
        Language);
  \item DialEgg~\cite{dialegg}, which translates any MLIR dialect to
        egglog and back.
\end{itemize}
Neither has been applied to the LAMP.
They are the natural alternative to a hand-written search
(\cref{sec:egraph}).

\paragraph{Library calls from MLIR.}
LAPIS~\cite{lapis} lowers MLIR linear algebra to Kokkos Kernels calls.
That is the ``op $\to$ library call'' step, but it does not use matrix
properties.

\paragraph{FLENS.}
FLENS~\cite{flens,flens-github} pursued the same goal, a convenient
notation with an efficient evaluation, as a header-only C++ library.
Overloaded operators build closures that are evaluated by BLAS calls: an
assignment such as \code{y = beta*y + alpha*transpose(A)*x + z} becomes
one \code{mv} (\code{gemv}) and one \code{axpy}, which inlining turns into
direct calls of the BLAS.
Matrix properties are part of the types (\code{GeMatrix},
\code{SyMatrix}, \code{TrMatrix}), and they select the kernel: \code{mv}
is \code{gemv}, \code{symv} or \code{trmv}, \code{mm} is \code{gemm},
\code{symm} or \code{trmm}.
By default FLENS never creates a temporary: an expression that BLAS cannot
evaluate without one is a compile-time error, or a run-time assertion where
the compiler cannot decide it.
LAPACK functionality is available through FLENS-LAPACK, called
explicitly (e.g.\ \code{lapack::sv} for a linear system).

FLENS also shows where a library has to stop.
Expression templates see one statement at a time, at the compile time of
C++, and matrix sizes only at run time.
They can map a fixed expression structure onto a fixed call, but they
cannot choose a parenthesisation by cost, replace an inverse by a
factorisation and share it, infer the properties of intermediate results,
find common subexpressions across statements, move invariant work out of
loops, or plan the memory of a whole function.
Those decisions need the whole program as data, the sizes and a cost
model: a compiler.
ulm.lam keeps the principles of FLENS and moves them there:
properties in types become \code{\#ulm_lam.props}; ``no hidden
temporaries'' becomes the caller-owned workspace with its query
(\cref{sec:workspace}) and the parameter intents that say what a function
overwrites; the notation becomes Python, traced at run time
(\cref{sec:python}); and the efficient evaluation is done by kernel
selection and ulmBLAS.

\paragraph{FLENS-LAPACK: the work of a compiler, by hand.}
FLENS-LAPACK re-implemented a significant part of LAPACK on top of
FLENS, bit for bit: compiled with \code{CHECK_CXXLAPACK}, every
FLENS-LAPACK function also calls the reference LAPACK on copies of its
arguments and requires identical results, compared by exact equality
(\file{flens/lapack/debug/isidentical.tcc}); in addition it is checked
with the Netlib LAPACK test suite (\file{flens/lapack/check-results.sh},
known deviations listed in \file{check-expected.txt}).
The work was so mechanical that it felt like doing the job of a
compiler: factorisations hoisted out of loops by hand, views of matrices
mapped onto workspace arrays by hand, and the size of every workspace
derived by hand through nested calls.
It also showed what LAPACK requires to avoid temporaries deep inside
nested calls: a workspace provided by the caller, a query that returns
its size (accumulated over the callees), in-place kernels on strided
views, and an \code{info} instead of an abort.
These requirements are the conventions of the code ulm.lam generates
(\cref{sec:workspace,sec:info}), and the mechanical work is what it
automates: loop-invariant code motion and the operator chosen by the
trip count (\cref{sec:kalman}), one factorisation per matrix, and the
workspace layout computed from the bufferized program.
So ulm.lam is the compiler that FLENS-LAPACK was written without.

% =====================================================================
\section{Architecture}
% =====================================================================

\begin{figure}[t]
\centering
\begin{tikzpicture}[
    node distance=5mm and 9mm,
    stage/.style={draw, rounded corners=2pt, fill=bg, align=left,
                  minimum width=52mm, inner sep=4pt, font=\small},
    pass/.style={font=\footnotesize\ttfamily, align=left},
    arr/.style={-{Latex[length=2mm]}, thick},
    note/.style={font=\footnotesize, align=left, text width=50mm}]
  \node[stage] (fe)  {frontend: domain-specific\\ language (DSL), Python bindings,\\or textual \code{ulm_lam} IR};
  \node[stage, below=of fe] (l1) {\code{ulm_lam} on tensors\\
     n-ary \code{mul}/\code{add}, \code{transpose}, \code{inverse}\\
     properties in \lstinline|#ulm_lam.props| encodings};
  \node[stage, below=12mm of l1] (k1) {\code{ulm_lamk} on tensors\\
     one op per kernel, DPS};
  \node[stage, below=12mm of k1] (k2) {\code{ulm_lamk} on memrefs};
  \node[stage, below=12mm of k2] (rt) {\code{emitc}: calls of\\
     \code{ulmBLAS::gemm}, \code{trlsm}, \code{potrf}, \dots};
  \node[stage, below=of rt] (ll) {C++17 $+$ ulmBLAS (header only)};

  \draw[arr] (fe) -- (l1);
  \draw[arr] (l1) -- node[pass, right=2mm] {ulm-lam-normalize, cse,\\
      ulm-lam-infer-properties,\\ ulm-lam-select-kernels} (k1);
  \draw[arr] (k1) -- node[pass, right=2mm] {ulm-lam-strip-properties,\\
      ulm-lam-results-to-out-params,\\ one-shot-bufferize,\\
      buffer-loop-hoisting,\\ ulm-lam-workspace} (k2);
  \draw[arr] (k2) -- node[pass, right=2mm] {convert-ulm-lamk-to-ulmblas} (rt);
  \draw[arr] (rt) -- node[pass, right=2mm] {mlir-translate --mlir-to-cpp} (ll);

  \node[note, left=8mm of l1] {algebra: normal form and property
    inference are \emph{local} rewrites (greedy is fine)};
  \node[note, left=8mm of k1.north west, anchor=north east, yshift=8mm]
    {the LAMP: \emph{global} search with a cost model; the only
     non-standard component};
  \node[note, left=8mm of k2] {memory: in-place updates and copies decided
    by upstream bufferization};
\end{tikzpicture}
\caption{The pipeline.  Everything below \code{ulm_lamk} on tensors is
upstream MLIR plus two small passes.}
\label{fig:pipeline}
\end{figure}

\Cref{fig:pipeline} shows the pipeline.
It splits into three concerns, and each uses the MLIR mechanism that fits
it:
\begin{description}[font=\normalfont\itshape]
  \item[Algebra] (\code{ulm_lam}).
        Normal form and property inference are local and confluent.
        Canonicalisation patterns and \code{InferTypeOpInterface} are the
        right tools.
  \item[Selection] (\code{ulm_lam} $\to$ \code{ulm_lamk}).
        A global optimisation problem.
        One pass owns it and keeps its own search data structure
        (\cref{sec:select}).
  \item[Memory and calls] (\code{ulm_lamk} $\to$ C++).
        Destination-passing style lets One-Shot Bufferize plan buffers.
        Each kernel becomes one call of the ulmBLAS C++ interface.
\end{description}

Two dialects, rather than one, make the boundary checkable.
After \code{ulm-lam-select-kernels}, a conversion target that marks
\code{ulm_lam} illegal proves that every operation was mapped.
Lowering passes never need to understand algebra.

% =====================================================================
\section{The \texttt{ulm\_lam} dialect}\label{sec:ulmlam}
% =====================================================================

\subsection{Values and properties}

Every operand is a rank-2 builtin tensor:
\begin{itemize}
  \item a column vector is $n\times1$ and a row vector is $1\times n$;
  \item scalars are plain floats;
  \item properties are a tensor \emph{encoding}:
\end{itemize}
\begin{lstlisting}
tensor<2000x2000xf64, #ulm_lam.props<spd>>
tensor<2000x200xf64,  #ulm_lam.props<full_rank>>
\end{lstlisting}
\code{sparse_tensor} attaches its storage format the same way.
The benefits:
\begin{itemize}
  \item standard tensor infrastructure (CSE, \code{tensor.empty},
        function signatures, the parser and printer) works unchanged;
  \item properties show up in the textual IR, so tests can
        \code{FileCheck} them.
\end{itemize}

The encoding is a bit set (\cref{lst:props}).
Its meaning follows three rules:
\begin{itemize}
  \item the set is an \emph{assertion}, and a missing bit means
        ``unknown'', never ``false''.
        Dropping bits is therefore always sound;
  \item sets are read modulo an implication closure
        (\code{ulm_lam::closure}), e.g.\ spd $\Rightarrow$ symmetric,
        non-singular;
  \item facts implied by a static shape (\code{square} on a
        $200\times200$ tensor) are not stored.
\end{itemize}
\code{full_rank} means $\operatorname{rank} = \min(\text{rows},\text{cols})$.
For square matrices that is \code{non_singular}.
The attribute implements \code{VerifiableTensorEncoding}, which rejects
contradictions such as \code{zero} together with \code{non_singular}.

\begin{figure}[t]
\srcbetween[style=td]{\src include/ulm_lam/UlmLamAttrs.td}{def UlmLam_PropNone}{// Kernel flags}{props}
\caption{The property encoding (\file{UlmLamAttrs.td}).}
\label{lst:props}
\end{figure}

\subsection{Property inference}\label{sec:props}

Every op implements \code{InferTypeOpInterface}.
Builders therefore compute result properties automatically, and every
newly created value is typed as strongly as local reasoning allows.
\Cref{tab:rules} shows the rules.
They are local, but they look at operand \emph{identity}, not only at
operand properties.
For example, $A\T A$ is symmetric positive semi-definite (SPSD) only because both
factors are the same SSA (static single assignment) value.

\begin{table}[t]
\centering\small
\begin{tabularx}{\linewidth}{@{}>{\raggedright}p{36mm}X@{}}
\toprule
Expression & Inferred properties (beyond closure) \\
\midrule
$A\T$ & lower\,$\leftrightarrow$\,upper swapped; all symmetric, definite,
  diagonal, orthogonal, rank properties preserved \\
$A\inv$ & requires non-singular; triangular, unit diagonal, diagonal,
  symmetric, spd, orthogonal, permutation preserved \\
$A_1\cdots A_k$ & lower (upper) if all are; unit diagonal, diagonal,
  orthogonal, permutation, non-singular if all are; zero if any is \\
$A\T A$ \;/\; $A A\T$ & spsd; spd if $A$ has full column (row) rank \\
$A\T P A$ & spd if $P$ spd and $A$ full column rank; spsd if $P$ spsd \\
$N A$, $A N$ & full rank if $A$ is and $N$ is non-singular \\
$A_1+\dots+A_k$ & symmetric / spsd / triangular / diagonal if all are;
  spd if all spsd and one spd \\
$\alpha A$ & structure preserved; spd, spsd only for constant $\alpha>0$;
  non-singular for constant $\alpha\neq0$ \\
\code{ulm_lamk} results & the same functions, through
  \code{PropertyInferenceOpInterface}: e.g.\ \code{trsm} result $L\inv X$ is
  full rank if $X$ is \\
\bottomrule
\end{tabularx}
\caption{Local property inference rules (selection).}
\label{tab:rules}
\end{table}

\paragraph{Relaxed type compatibility.}
By default, \code{InferTypeOpInterface} requires the result type to
\emph{equal} the inferred type.
That rule would make rewriting fragile.
Replacing an operand by a value with stronger properties would invalidate
every user until its type was recomputed.
The \code{ulm_lam} ops therefore override \code{isCompatibleReturnTypes}.
A result type is accepted if its shape matches and its properties are
\emph{implied by} the inferred ones (\cref{lst:mul}).
Weaker types are always valid.
\code{ulm-lam-infer-properties} strengthens all types in one topological
sweep.
Facts that inference cannot derive enter only through function signatures
or \code{ulm_lam.assume}.
An assumption is also tested at run time where it matters: a Cholesky
factorisation of a matrix that is assumed spd reports through \code{info}
when the matrix is not positive definite numerically
(\cref{sec:info}).

\subsection{Operations}

\begin{table}[t]
\centering\small
\begin{tabularx}{\linewidth}{@{}lX@{}}
\toprule
Op & Semantics \\
\midrule
\code{ulm_lam.mul} & n-ary product $A_1\cdots A_k$; associative, no parenthesisation \\
\code{ulm_lam.add} & n-ary sum; associative and commutative (\code{Commutative} trait) \\
\code{ulm_lam.scale} & $\alpha A$; subtraction is \code{add(A, scale(-1, B))} \\
\code{ulm_lam.transpose} & $A\T$ \\
\code{ulm_lam.inverse} & $A\inv$; verifier requires square; never computed unless requested \\
\code{ulm_lam.identity} & identity matrix (with dynamic sizes as operands) \\
\code{ulm_lam.assume} & add properties (trusted, or \code{checked} at run time) \\
\bottomrule
\end{tabularx}
\caption{\code{ulm_lam} operations. All are \code{Pure}.}
\label{tab:ulmlamops}
\end{table}

\Cref{tab:ulmlamops} lists the operations.
\code{ulm_lam.mul} is n-ary on purpose.
A product denotes the mathematical object $A_1\cdots A_k$.
Choosing a parenthesisation is a decision of kernel selection, not of the
frontend.

\begin{figure}[t]
\srcdef{\src include/ulm_lam/UlmLamOps.td}{UlmLam_MulOp}
\caption{\code{ulm_lam.mul} in ODS.  The common base class (not shown)
adds \code{Pure}, \code{InferTypeOpInterface} with relaxed compatibility,
and \code{PropertyInferenceOpInterface}.}
\label{lst:mul}
\end{figure}

\subsection{Normal form}\label{sec:normal}

\code{ulm-lam-normalize} applies canonicalisation patterns greedily and then
runs CSE.
The patterns are:
\begin{enumerate}
  \item flatten nested \code{mul} and \code{add};
  \item push transposes to the leaves:
        $(AB)\T\to B\T A\T$, $(A+B)\T\to A\T+B\T$, $(A\T)\T\to A$, and
        $A\T\to A$ for symmetric $A$;
  \item push inverses to the leaves for square non-singular factors:
        $(AB)\inv\to B\inv A\inv$ and $(A\T)\inv\to(A\inv)\T$.
        A leaf is then $A$, $A\T$, $A\inv$ or $A\invT$;
  \item hoist scales out of products and merge them;
  \item remove identities from products and zeros from sums;
        a product with a zero factor is zero.
\end{enumerate}
Each pattern strictly decreases a well-founded measure: the depth of
transpose and inverse nodes, and the nesting of \code{mul} and \code{add}.
The driver therefore terminates.
Normalising is a precondition for inference.
In \cref{sec:example}, the frontend's binary product $(X\T M\inv)X$ is not
recognised as spd.
The flattened product $X\T M\inv X$ is.

The normal form is not unique for the purposes of the LAMP.
Distributing a product over a sum can help or hurt, and so can factoring
one out.
Such rewrites are therefore \emph{not} canonicalisations.
They are moves of the search.

% =====================================================================
\section{The \texttt{ulm\_lamk} dialect}\label{sec:ulm_lamk}
% =====================================================================

\code{ulm_lamk} has one operation per kernel (\cref{tab:ulm_lamk}), with BLAS
flags as enum attributes.
The design follows four rules.

\paragraph{Destination-passing style.}
The operand that BLAS overwrites is the \code{outs} operand: $C$ for
\code{gemm}, $B$ for \code{trsm}, $A$ for \code{potrf}.
On tensors, the op returns the updated value.
On memrefs, it returns nothing.
This is the contract of \code{DestinationStyleOpInterface}.
It is exactly what One-Shot Bufferize needs to update \code{outs} buffers
in place and to copy only when an overwritten value is still live
(\cref{sec:lowering}).

\paragraph{Precise partial results.}
\code{potrf} and \code{syrk} define only one triangle.
On tensors, their semantics are those of BLAS: the \code{uplo} triangle
receives the result, and the strict other triangle keeps the init's values.
With a \code{tensor.empty} init those values are undefined.
Kernel selection guarantees that every consumer reads only the defined
triangle: the factor of \code{potrf} is only used by triangular solves
with the same \code{uplo}, and the result of \code{syrk} is passed on as a
matrix stored in its lower triangle (next paragraph).

\paragraph{Symmetric matrices stored as one triangle.}\label{sec:storage}
By default a matrix is stored completely, also a symmetric one, and the
generated code may read every element (e.g.\ with \code{gemm}).
As with the \code{uplo} argument of LAPACK, a parameter can instead be
declared to be stored as one triangle: \code{Matrix(n, n, "spd",
stored="lower")}, in the encoding \code{\#ulm_lam.props<spd, lower_stored>}.
The other triangle may then hold anything.
Unlike the other bits of the encoding, these two are not assertions about
the value but about its memory, so dropping them is \emph{not} sound.
They are never inferred and never propagated to the result of an op;
kernel selection resolves every use:
\begin{itemize}
  \item a product with such a factor uses \code{symm} (ulmBLAS
        \code{sylmm}/\code{syumm}; the right side $BA$ is the left one on
        transposed views, $C\T = AB\T$) or \code{symv}
        (\code{sylmv}; for the upper triangle on the transposed view,
        because $A\T = A$), never \code{gemm}/\code{gemv};
  \item a Cholesky factorisation uses the stored triangle:
        $A = LL\T$, or $A = U\T U$ and $A\inv = U\inv U\invT$;
  \item where a kernel reads the whole matrix (a product of two such
        factors, sums, right-hand sides of solves, LU, and uses outside
        the selected code: returns, calls, loops) the matrix is completed
        by \code{ulm_lamk.sycomplete}, which mirrors the stored triangle.
        It is created once, right after the definition of the matrix, and
        shared by all such uses.  As every DPS op it works in place, and
        One-Shot Bufferize decides where: on the matrix itself if it may
        be overwritten (a temporary, or an \code{inout} parameter) and is
        not read afterwards, on a copy otherwise.
        Of two such factors of one product, the one that is complete
        already, else one that can be completed in place, is completed.
\end{itemize}
The interface comment names the triangle (\code{read only, lower triangle
(symmetric)}).
The tests fill the other triangle with NaN, so that any kernel reading it
shows in the result.

The result of \code{syrk} is handled the same way.
A Gram product $Z\T Z$ that makes up a whole product and goes into a new
matrix is computed by \code{syrk} ($n^2k$ instead of $2n^2k$ FLOPs, and
the DP prices it so); its result is passed on as stored in the lower
triangle.  The following Cholesky factorisation reads that triangle
anyway; a return or a sum gets it completed (an $n^2$ copy instead of
$n^2k$ FLOPs).
Gram products inside a larger chain, or reused by later products, stay
\code{gemm}.

\paragraph{Costs.}
Every op implements \code{KernelCostOpInterface}, with FLOPs and bytes
moved.
Sizes come from static shapes.
For dynamic dimensions (\cref{sec:dynamic}) the costs depend on run-time
sizes, and kernel selection chooses at run time (next paragraph).
Multi-versioning with a run-time dispatch on sizes is future work.

\paragraph{Scalars.}
\code{alpha} and \code{beta} are SSA values whose type must match the
matrix element type (\code{TypesMatchWith}).
Scales and accumulations folded into kernels need no special case.

\begin{table}[t]
\centering\small
\begin{tabularx}{\linewidth}{@{}llX>{\raggedleft\arraybackslash}p{30mm}@{}}
\toprule
Op & ulmBLAS & Semantics & FLOPs \\
\midrule
\code{gemm} & \code{gemm} & $C \leftarrow \alpha\,\mathrm{op}(A)\,\mathrm{op}(B)+\beta C$ & $2mnk$ \\
\code{symm} & \code{sylmm} & $C \leftarrow \alpha S B+\beta C$, $S$ symmetric & $2m^2n$ \\
\code{trmm} & \code{trlmm}/\code{trumm} & $B \leftarrow \alpha\,\mathrm{op}(L) B$ & $m^2n$ \\
\code{trsm} & \code{trlsm}/\code{trusm} & $B \leftarrow \alpha\,\mathrm{op}(L)\inv B$ & $m^2n$ \\
\code{syrk} & \code{sylrk} & $C \leftarrow \alpha A\T A+\beta C$ (one triangle) & $n^2k$ \\
\code{gemv} & \code{gemv} & $y \leftarrow \alpha\,\mathrm{op}(A)x+\beta y$ & $2mn$ \\
\code{trsv} & \code{trlsv}/\code{trusv} & $x \leftarrow \mathrm{op}(L)\inv x$ & $n^2$ \\
\code{potrf} & \code{potrf} & $A = LL\T$, in place & $n^3/3$ \\
\code{getrf} & \code{getrf} & $PA = LU$ (lower--upper), in place & $2n^3/3$ \\
\code{getrs} & \code{getrs} & $B \leftarrow \mathrm{op}(A)\inv B$ from LU & $2n^2k$ \\
\code{geaxpy} & \code{geaxpy} & $Y \leftarrow \alpha X + Y$ (matrices) & $2mn$ \\
\code{gescal} & \code{gescal} & $X \leftarrow \alpha X$ (matrices) & $mn$ \\
\code{call} & -- & external function updating its \code{outs} in place & -- \\
\bottomrule
\end{tabularx}
\caption{\code{ulm_lamk} operations and the ulmBLAS routines they are
lowered to.  The BLAS-like operations keep the familiar \code{trans} and
\code{side} attributes; in the lowering they become strides
(\cref{sec:ulmblas}).  \code{trmm} is defined but not yet produced by
kernel selection.  Not listed: \code{symv} (\code{sylmv}) and
\code{sycomplete} (a helper in the generated code), see
\cref{sec:storage}.}
\label{tab:ulm_lamk}
\end{table}

\begin{figure}[t]
\srcdef{\src include/ulm_lam/UlmLamKernelOps.td}{UlmLamk_TrsmOp}
\caption{\code{ulm_lamk.trsm} in ODS.  The base class \code{UlmLamk_Op} adds
\code{DestinationStyleOpInterface} (with \code{B} as init), memory
effects, property inference and the cost interface.}
\label{lst:trsm}
\end{figure}

% =====================================================================
\section{Kernel selection}\label{sec:select}
% =====================================================================

\subsection{Problem}

The input is a function of \code{ulm_lam} ops in normal form, with
properties inferred.
The output is a sequence of \code{ulm_lamk} ops (and \code{tensor.empty})
that computes the same values at minimal cost.
Plain pattern rewriting cannot do this for three reasons:
\begin{itemize}
  \item kernels match \emph{subsequences} of n-ary products;
  \item choices interact.
        A Cholesky factorisation of $M$ pays off only because both
        occurrences of $M\inv$ reuse it;
  \item the best choice is visible only in the total cost.
\end{itemize}

\subsection{The kernel catalogue}

Kernels are declared, not hand-coded (\cref{lst:catalogue}).
A record has five parts:
\begin{itemize}
  \item a pattern over the normal form;
  \item constraints on properties, vector shapes and operand identity;
  \item the \code{ulm_lamk} op and its flags;
  \item the operand that is overwritten;
  \item a FLOP formula.
\end{itemize}
A custom TableGen backend (\code{ulm-lam-tblgen -gen-kernel-catalogue}) would
generate the matchers, cost functions and builders.
Today the file is checked for syntax with \code{mlir-tblgen
--print-records}, which yields 31 records.
Supporting another library, or a fused vendor kernel, means adding
records.
Linnea's kernel description plays the same role.

A \code{mul} pattern with $k$ operands matches any $k$ \emph{consecutive}
factors of a longer product.
This restricted form of associative matching is all that products need,
and it is linear in the product length.
Commutative matching is needed only for sums.
Sums are short in practice, so subsets can be enumerated.
Full associative-commutative matching in the style of MatchPy is not
required.

Factorisations are expression rewrites, as in Linnea.
Cholesky replaces an spd leaf $A$ by $LL\T$, with $L$ lower triangular and
non-singular.
Normal form then turns $A\inv$ into $L\invT L\inv$.
The ordinary \code{trsm}/\code{trsv} patterns pick up the solves, and no
kernel needs to know about factorisations.

\begin{figure}[t]
\srcbetween[style=td]{\src include/ulm_lam/UlmLamKernels.td}{// Triangular solves with a matrix}{// BLAS 2}{trsm-cat}
\srcdef{\src include/ulm_lam/UlmLamKernels.td}{Cholesky}
\caption{Excerpt of the kernel catalogue (\file{UlmLamKernels.td}).}
\label{lst:catalogue}
\end{figure}

\subsection{Inner solver: generalised matrix chain}

For a single product $A_1\cdots A_k$, whose leaves are decorated with
$\T$ and $\inv$, the GMC algorithm~\cite{gmc} is a dynamic programme over
intervals:
\[
  \mathrm{cost}(i,j) \;=\; \min_{i\le s<j}\;
  \min_{\kappa \in \mathcal{K}(T_{i,s},\,T_{s+1,j})}
  \bigl[\mathrm{cost}(i,s)+\mathrm{cost}(s+1,j)+\mathrm{flops}_\kappa\bigr],
\]
where:
\begin{itemize}
  \item $T_{i,s}$ is the temporary for the interval $A_i\cdots A_s$,
        with the properties inferred for it;
  \item $\mathcal{K}$ is the set of catalogue kernels that match the
        binary product of the two temporaries.
\end{itemize}
Because temporaries carry inferred properties, structure propagates
through the chain.
For example, $L\inv X$ stays full rank, so $(L\inv X)\T(L\inv X)$ becomes
\code{syrk} with an spd result.
The cost is $O(k^3\,|\mathcal{K}|)$.
GMC alone is a complete but myopic strategy (\code{strategy=gmc}).

\subsection{Implemented: \texttt{strategy=gmc}}\label{sec:impl}

The implementation of \code{ulm-lam-select-kernels}
(\file{lib/UlmLamPasses.cpp}) handles products, sums and inverses:
\begin{itemize}
  \item \textbf{Inverses.}
        Every inverted matrix $M$ is factorised once, right after its
        definition, chosen by its properties.
        A triangular $M$ is used as it is; for an spd $M$, \code{potrf}
        gives $M=LL\T$ and $M\inv$ becomes the two solves $L\invT L\inv$;
        any other $M$ is factorised by \code{getrf}, and $M\inv$ becomes an
        LU solve.
        The factorisation is shared by all products in which $M\inv$
        occurs, and it stays outside a loop in which $M$ is invariant.
        An inverse is never formed explicitly; an inverse that is not a
        factor of a product is diagnosed.
  \item \textbf{Products.}
        Every product is parenthesised by the generalised matrix chain
        dynamic programme above.
        An interval is computed by a \code{gemm} (\code{gemv} for a
        vector), by a solve applied from the left to the rest (\code{trsm},
        \code{trsv}, \code{getrs}, in place), by a triangular solve applied
        from the right (\code{trsm} with \code{side(right)}), or as $Z\T Z$
        when its left half is the transpose of its right half.
        The last case is Linnea's common subexpression up to
        transposition: $X\T L\invT L\inv X = (L\inv X)\T(L\inv X)$.
        Transposition is always a kernel flag, never a copy.
        For example, $A\,B\,x$ with $1000\times1000$ matrices becomes
        $A(Bx)$, and $A\inv B C$ becomes $A\inv(BC)$ when $C$ has fewer
        columns than $B$.
  \item \textbf{Reuse across products.}
        Every interval that a product computes as a plain product is
        remembered.
        In a later product, the same interval, or its transpose, costs
        nothing if the remembered value dominates the product.
        So the second product of GLS reuses $W=L\inv X$ from the first:
        $X\T L\invT L\inv y = W\T(L\inv y)$, one triangular solve
        instead of two.
        This is a first step towards the sharing that Linnea's search
        finds (\cref{sec:search}).
  \item \textbf{Sums.}
        $Y_1+a_2Y_2+\dots+b_1P_1+\dots$ accumulates into the first plain
        term: \code{gescal} for its scale, \code{geaxpy} for every further
        plain term, and the last kernel of every product writes into the
        accumulator (\code{gemm} with $\beta=1$; a solve through
        \code{geaxpy}).
        $X+2A\T B$ is therefore a single \code{gemm} with
        \code{trans_a(trans)}, $\alpha=2$ and $\beta=1$ that writes into
        $X$.
  \item \textbf{Diagnosed, not miscompiled:} explicit inverses, right
        solves with an LU factorisation, transposed accumulators, and
        transposes that are not kernel operands.
\end{itemize}
$Z\T Z$ uses \code{syrk} where its result goes into a new matrix
(\cref{sec:storage}).
An inverse of a product of square, non-singular factors,
$(F_1\cdots F_k)\inv$, is inverted factor by factor,
$F_k\inv\cdots F_1\inv$, when factorising the distinct factors costs less
than forming the product and factorising it: $(RR)\inv$ with $R$ spd is one
Cholesky factorisation of $R$ ($n^3/3$ instead of $2n^3 + n^3/3$, and
$\kappa(R)$ instead of $\kappa(R)^2$).
Products that read the same backwards when transposed,
$F_1\cdots F_k$ with $F_i\T = F_{k+1-i}$, are symmetric (if the middle
factor is), spsd, or spd: $A\T A$, $A\T PA$, but also $RR$ and $R^3$.

The other passes are \code{ulm-lam-normalize} (flattening of products and
sums, folding of double transposes and inverses, of transposes of
symmetric matrices and of scaling by one), \code{ulm-lam-infer-properties}
and \code{ulm-lam-strip-properties}.
The dialects implement the property inference rules of \cref{tab:rules}
for products, sums, scaling, transposes and inverses, the relaxed type
compatibility, the verifiers, FLOP and byte costs for every kernel, and
bufferization of every \code{ulm_lamk} op in destination-passing style.

\subsection{Outer solver: search over derivations}\label{sec:search}

The default strategy is Linnea's search, ported.
\begin{itemize}
  \item \textbf{State}: a hash-consed expression DAG of what remains to
        compute, plus the kernels emitted so far and their cost.
        Hash-consing makes CSE free.
        Nodes are canonicalised \emph{modulo transposition}, so
        $X\T L\invT$ is found as $(L\inv X)\T$ once $W=L\inv X$ exists.
  \item \textbf{Moves}:
        \begin{enumerate*}[label=(\roman*)]
          \item apply a factorisation to a leaf
                (all its occurrences change together);
          \item compute a subexpression with GMC and replace it by a new
                leaf whose properties are inferred;
          \item apply a non-canonical algebraic rewrite
                (distribute, factor out).
        \end{enumerate*}
  \item \textbf{Order}: best-first by cost so far plus a lower bound on
        the cost of the remainder.
        The search stops on a time budget (\code{budget-ms}).
        Like Linnea, it is anytime: the first complete derivation comes
        from a greedy descent, and better ones replace it.
\end{itemize}

The search works on its own data structure, not on IR.
Cloning functions for every state would be wasteful.
A search state is also not a valid program: partly selected expressions
mix factors and pending products.
The DAG reuses the C++ property functions of the ops, so inference in the
search and in the IR cannot disagree.
Only the winning derivation is materialised as \code{ulm_lamk} ops.
A debug option dumps the explored derivation graph.

\subsection{Alternative: equality saturation}\label{sec:egraph}

\code{strategy=egraph} is the experimental route.
It translates the \code{ulm_lam} function into an e-graph, through the
\code{eqsat} dialect~\cite{eqsat} or DialEgg/egglog~\cite{dialegg}.
It saturates with the algebraic identities and the catalogue rules, then
extracts the cheapest term.
The idea is attractive, because the moves of \cref{sec:search} become
plain rewrite rules.
There are three difficulties:
\begin{itemize}
  \item properties are per e-class facts.
        egglog lattices with \code{:merge} can represent them, but
        rewrites become conditional on analyses;
  \item extraction with sharing is NP-hard, and integer programming is
        the usual approach.
        This matches the NP-completeness of the LAMP;
  \item associativity makes the e-graphs of products blow up.
        GMC as a custom extractor for chain e-classes is one way out.
\end{itemize}
A comparison of both strategies on the LAMP benchmark~\cite{lamp-paper} would
itself be a contribution.

\subsection{Soundness and numerics}

Kernel selection must not change the result beyond rounding.
Rewrites with side conditions are checked against inferred properties:
\begin{itemize}
  \item $(AB)\inv = B\inv A\inv$ only for square non-singular factors;
  \item $A\T=A$ only for symmetric $A$.
\end{itemize}
Explicit inverses are never formed unless the inverse itself is an
output.
Factorisations are chosen by structure: Cholesky for spd, LU with pivoting
otherwise, and QR (orthogonal $Q$ times upper triangular $R$) where
requested.
Every factorisation reports failure through the \code{info} of the
function (\cref{sec:info}), so a violated property assumption is
reported to the caller instead of producing garbage.

% =====================================================================
\section{Bufferization and lowering}\label{sec:lowering}
% =====================================================================

\paragraph{Buffers.}
After \code{ulm-lam-strip-properties}, the IR consists of \code{ulm_lamk} ops on
plain tensors.
\code{ulm-lam-results-to-out-params} first makes every matrix result a
by-reference parameter (\cref{sec:loops}).
The pipeline then runs upstream passes in four steps:
\begin{enumerate}
  \item \code{one-shot-bufferize\{bufferize-function-boundaries\}};
  \item \code{buffer-results-to-out-params\{hoist-static-allocs\}};
  \item \code{buffer-loop-hoisting}, which moves allocations of
        loop-invariant size out of loops (\cref{sec:loops});
  \item \code{buffer-deallocation-pipeline}.
\end{enumerate}
\code{ulm-lam-results-to-out-params} marks the matrices passed by value
\code{bufferization.writable = false}.
The analysis therefore copies an input only when a kernel overwrites it.
In the example these are $M$ (\code{potrf}), $X$ (\code{trsm}) and $y$
(\code{trsv}).
Intermediate results that end in the output are computed directly in the
result buffer.
That is Linnea's memory optimisation, here obtained from upstream
bufferization.

\subsection{Lowering to ulmBLAS}\label{sec:ulmblas}

\code{convert-ulm-lamk-to-ulmblas} (\file{lib/ConvertUlmLamkToUlmBLAS.cpp})
lowers the bufferized program to EmitC, MLIR's dialect for emitting C and
C++.
Every kernel becomes one call of the C++ interface of
ulmBLAS~\cite{ulmblas}, and \code{mlir-translate --mlir-to-cpp} prints
C++17.
Loops, arithmetic and functions use the upstream
\code{\{scf,arith,func\}-to-emitc} patterns with the type converter of the
pass.

\paragraph{Why ulmBLAS and not CBLAS/LAPACK.}
ulmBLAS is header only, its element type is a template parameter, and
every matrix is a pointer with two strides (\code{incRow}, \code{incCol}),
for BLAS and LAPACK routines alike.
Strides are what an MLIR \code{memref} carries, so the lowering is direct:
\begin{itemize}
  \item a transposed operand is the same pointer with swapped strides, so
        \code{trans} flags disappear, and the transpose of a lower
        triangular matrix is an upper triangular one
        (\code{trusv} on $L$ with strides $(1,n)$ solves with $L\T$);
  \item a triangular solve from the right is a solve from the left on
        transposed views: $B\,\mathrm{op}(T)\inv =
        (\mathrm{op}(T)\invT B\T)\T$, so \code{side} disappears too;
  \item row-major versus column-major is not an issue: LAPACK routines
        such as \code{potrf}, \code{getrf} and \code{getrs} take strides
        as well.
        With CLAPACK, row-major data needs tricks (a row-major lower
        triangle is the column-major upper one; \code{getrf} factorises
        $A\T$) and \code{getrs} with several right-hand sides cannot be
        used at all;
  \item factorisations return \code{info}, which the generated function
        turns into its own \code{info} (\cref{sec:info}).
\end{itemize}

\paragraph{A second lowering: CBLAS and LAPACK.}\label{sec:cblas}
ulm.lam does two things at once: it maps formulas onto BLAS and LAPACK
routines (the LAMP), and it runs them on ulmBLAS.
To measure the first independently of the second, the option
\code{backend=cblas} of \code{convert-ulm-lamk-to-ulmblas} (\code{BACKEND=cblas}
for \file{scripts/ulm-lam-compile.sh}, \code{lam.compile(...,
backend="cblas")} in Python) emits the same calls with the namespace
\code{ulm\_lam\_cblas} instead of \code{ulmBLAS}.
The header \file{include/ulm\_lam/cblas\_backend.h} implements that
interface with CBLAS and the Fortran LAPACK symbols, which MKL (Intel Math
Kernel Library), OpenBLAS and Accelerate all provide.
It works in column-major order and uses the tricks above instead of
copies: a row-major operand is the transposed column-major one (a trans
flag, a flipped triangle, or the transposed problem $C\T = B\T A\T$),
\code{potrf} factorises the flipped triangle, and \code{getrf} factorises
$A\T$, which \code{getrs}, called with the same strides, takes into
account.
Only a right-hand side of \code{getrs} and the $B$ of \code{symm} whose
orientation differs from that of $C$ are copied ($O(nm)$).
The kernel sequence is identical, so the time of the generated code on
MKL compared with the hand-written MKL version measures the mapping, and
the same code on ulmBLAS compared with it on MKL measures ulmBLAS
(\cref{sec:gpcluster}).
All 44 Python tests pass with both backends.

\paragraph{The ulm.lam ABI.}
\begin{itemize}
  \item A static row-major \code{memref} is a pointer (\code{double *},
        \code{float *}, pivots \code{int32_t *}) with strides
        (columns, 1); \code{index} is \code{size_t}.
        So \code{foo(index, memref<500x500xf64>)} is
        \code{void foo(size_t, double *)}.
  \item Matrices are passed by reference.
        A function with a matrix output receives a pointer to it; the
        caller allocates and frees the matrix and owns it.
        No function returns a matrix, so ownership never moves from callee
        to caller, and an external function cannot leak.
        \code{ulm-lam-results-to-out-params} and
        \code{buffer-results-to-out-params} establish the convention;
        \code{convert-ulm-lamk-to-ulmblas} rejects a function that violates
        it.
  \item Every matrix parameter has an intent, \code{in}, \code{inout}
        or \code{out}, declared in the parameter list (see below).
  \item An external function that updates a matrix receives it by
        reference and changes it in place: a declaration whose result is
        marked \code{\{ulm_lam.inout = k\}} becomes a \code{ulm_lamk.call}
        (\cref{sec:adaptive}).
  \item A function that needs memory for temporaries gets it from the
        caller, in the style of LAPACK's \code{lwork}
        (\cref{sec:workspace}).
  \item Functions are \code{extern "C"}: C code calls the generated C++,
        and the generated C++ calls C functions.
\end{itemize}
\paragraph{Parameter intents.}
LAPACK routines commonly overwrite the matrices they receive; the caller
learns it from the documentation (``on exit, A is destroyed'').
In ulm.lam the parameter list says it:
\begin{itemize}
  \item \code{\{ulm_lam.in\}} (the default for matrix arguments): read only.
        If a kernel would overwrite the matrix, the function works on a
        copy, which lives in its workspace.
  \item \code{\{ulm_lam.inout\}}: the function may overwrite the matrix and
        use it as workspace.
        If a result is marked \code{\{ulm_lam.inout = k\}}, argument $k$ holds
        that result on exit (the loop example updates $X$ this way);
        otherwise its contents on exit are unspecified.
  \item \code{\{ulm_lam.out\}} (the default for results): an out parameter;
        the caller provides the buffer, the function writes it and also
        computes in it (\code{gemv} and the triangular solves of GLS work
        in the caller's $b$).
\end{itemize}
\code{ulm-lam-results-to-out-params} turns the intents into bufferization
attributes (\code{in} arguments are not writable, \code{inout}
arguments are), so One-Shot Bufferize decides copies and in-place updates
accordingly.
\file{examples/gls/gls_inout.mlir} is GLS with $M$ and $y$ declared
\code{inout}: the Cholesky factorisation of $M$ and the solve with $y$
work on the caller's matrices, only $X$ (\code{in}) is copied for the
in-place solve $W=L\inv X$, and the workspace shrinks from 35.5\,MB to
3.5\,MB.
The result is bit-identical to that of \code{gls}.

Because the caller has to know the intents, the generated code states
them: every public function is preceded by a comment generated from its
parameter attributes,
\begin{lstlisting}[style=c]
/* gls_inout: matrices by reference (row-major), owned by the caller
 *   v1  double[2000][200]   in     read only
 *   v2  double[2000][2000]  inout  may be overwritten (unspecified on exit)
 *   v3  double[2000][1]     inout  may be overwritten (unspecified on exit)
 *   v4  double[200][1]      out    written
 *   v5  int8_t *            work   workspace (bytes)
 *   v6  int64_t *           info   -1: query, 0: allocate, >0: size of work
 */
\end{lstlisting}

Memory: \code{memref.alloc}, \code{dealloc} and \code{copy} become
\code{aligned_alloc} (size rounded up to the alignment, as C++17
requires), \code{free} and \code{ulm_lam_copy}, a \code{memcpy} that
copies large arrays in parallel when compiled with OpenMP.
On a machine with several NUMA (Non-Uniform Memory Access) domains the
first write decides where a page of a new buffer lies: a serial copy
into a workspace the function allocated itself would put all of it on
one domain, and the kernels that then use it with all threads pay for
that (\cref{sec:gpcluster}).
Upstream \code{convert-memref-to-emitc} is not used: it types arguments as
C arrays (\code{double[500][500]}) but allocations as pointers, which
leaves unresolved casts.

\paragraph{Dynamic sizes.}\label{sec:dynamic}
Matrix sizes can be run-time parameters.  The argument (and result)
attribute \code{ulm_lam.dims} ties every dynamic dimension of a matrix to
an index parameter of its function (argument positions, one per \code{?}):
\begin{lstlisting}[style=mlir]
func.func @gls(%n: index, %m: index,
               %X: tensor<?x?xf64, #fr> {ulm_lam.dims = [0, 1]},
               %M: tensor<?x?xf64, #spd> {ulm_lam.dims = [0, 0]},
               %y: tensor<?x1xf64> {ulm_lam.dims = [0]})
    -> (tensor<?x1xf64> {ulm_lam.dims = [1]})
\end{lstlisting}
The generated function takes the sizes as \code{size_t} parameters,
\code{gls(size_t v1, size_t v2, double *v3, ...)}, and the interface
comment says \code{double[v1][v2]}.
\code{ulm_lam::dimOf} traces the size of a dimension through the ops of
the pipeline to a parameter, an allocation, or a call result (whose
callee names the size by \code{ulm_lam.dims} on its result); kernel
selection, the workspace layout and the lowering use it.
The results of a public function become out parameters before
bufferization, and with \code{eliminate-empty-tensors} the kernels write
into the caller's buffer directly (\code{buffer-results-to-out-params}
can only do that for static shapes).
The workspace query is computed from the sizes and returns what exactly
the evaluation chosen by the compiler needs.
All ulmBLAS calls name their index type, \code{ulmBLAS::gemm<std::ptrdiff_t>},
so that literals and \code{size_t} parameters convert to it.
\file{examples/dynamic} checks one compiled GLS for six sizes, from
$1\times1$ to $2000\times200$.

\paragraph{Plans chosen at run time.}
Which evaluation is best can depend on the sizes: for $D = ABC$ with
$A\in\mathbb{R}^{n\times k}$, $B\in\mathbb{R}^{k\times m}$, $C\in\mathbb{R}^{m\times p}$,
$(AB)C$ costs $2nkm+2nmp$ and $A(BC)$ costs $2kmp+2nkp$.
Kernel selection runs the matrix-chain DP for every assignment of
$\{10, 1000\}$ to the (at most four) distinct dynamic sizes of a product;
if the optimal plans differ, the generated code computes the cost of each
from the actual sizes (index arithmetic, as the DP does) and runs the
cheapest (an \code{scf.if} chain):
\begin{lstlisting}[style=c]
if (((v1 * v2) * v3) * 2 + ((v1 * v3) * v4) * 2 <=
    ((v2 * v3) * v4) * 2 + ((v1 * v2) * v4) * 2) {
    ulmBLAS::gemm<std::ptrdiff_t>(v1, v3, v2, ..., v5, ..., v6, ..., T, ...);  // T = A B
    ulmBLAS::gemm<std::ptrdiff_t>(v1, v4, v3, ..., T, ..., v7, ..., v8, ...);  // D = T C
} else {
    ulmBLAS::gemm<std::ptrdiff_t>(v2, v4, v3, ..., v6, ..., v7, ..., T, ...);  // T = B C
    ulmBLAS::gemm<std::ptrdiff_t>(v1, v4, v2, ..., v5, ..., T, ..., v8, ...);  // D = A T
}
\end{lstlisting}
Every plan writes the product into the caller's buffer.
The workspace parts of the plans overlap, and the workspace query evaluates
the same condition, so it returns the size of exactly the plan that runs:
\code{v16 ? round64(8 * v1 * v3) : round64(8 * v2 * v4)}.
Within a plan no subproduct is reused by later products (it would not
dominate them), and no \code{syrk} result is passed on as one triangle.
\file{examples/dynamic/driver\_chain.cc} checks sizes for which either plan
wins.

\paragraph{Workspaces.}\label{sec:workspace}
A generated function does not allocate memory behind the caller's back.
\code{ulm-lam-workspace} (\file{lib/UlmLamWorkspace.cpp}) runs after
bufferization.
A public function with temporaries gets two trailing parameters:
\begin{lstlisting}[style=c]
void gls(double *X, double *M, double *y, double *b,
         int8_t *work, int64_t *workspace_info);
\end{lstlisting}
\begin{itemize}
  \item \code{*workspace_info == -1}: a query.
        The function sets \code{*workspace_info} to the required size in
        bytes and returns without computing anything.
  \item \code{*workspace_info == 0}: the function allocates one block of
        the required size, uses it, and frees it before it returns.
  \item \code{*workspace_info > 0}: the size of \code{work}.
        It must be at least the required size; if it is too small, the
        function returns $\mathit{info} = -i$, where $i$ is the position of
        \code{workspace_info} (as LAPACK reports an illegal
        \code{lwork}), and computes nothing.
        Otherwise all temporaries live in \code{work}.
\end{itemize}
The sizes are in bytes, because the temporaries mix element types
(matrices and \code{int32_t} pivots).
All temporaries, including those of loops, become slots of the one block
at 64-byte aligned offsets (computed at compile time, or from the sizes
for dynamic dimensions); in the C++ each is
\code{(double *)(work + offset)}.
The query returns the size that exactly the evaluation chosen by the
compiler needs; there is no separate minimal size as in LAPACK.
Functions without temporaries keep their plain signature.

\paragraph{Accumulation.}
As a LAPACK routine that calls other routines, a function that calls
functions with a workspace answers a query with its own requirement plus
those of its callees.
\code{ulm-lam-workspace} decides transitively which functions take a
workspace: those with temporaries, and those that call a function with a
workspace.
An external function declares that it follows the convention with the
attribute \code{\{ulm_lam.workspace\}}; its last two parameters are then
\code{work} and \code{workspace_info}.
At entry, a function queries each such call (\code{*workspace_info = -1};
the matrix arguments are dummies, which a query does not touch; scalars
that are not available at entry, such as a loop index, are 0) and
restores its own \code{*workspace_info}.
Its block is laid out as [own temporaries | callee 1 | callee 2 | \dots],
each part 64-byte aligned, and every call gets its part and its size.
Afterwards the function restores \code{*workspace_info}, as LAPACK leaves
\code{lwork} unchanged.
So the requirement of an external function with a workspace may depend
only on its matrix sizes, which are static in ulm.lam, and on scalar
arguments that are available at entry.
Recursion between functions with a workspace is rejected.

\file{examples/compose/compose.mlir} has a generated function
\code{estimate_all} that calls the generated \code{gls} and the external
\code{postprocess} in a loop (\cref{lst:compose}).
Its query returns 35\,555\,200 bytes: its own 17\,600 ($y_i$ and $b_i$)
plus the 35\,536\,000 of \code{gls} plus the 1\,600 of
\code{postprocess}.
\file{examples/compose/driver.cc} checks the accumulation, that every
call of \code{postprocess} gets a large enough part of the caller's
workspace, that the estimates equal separate calls of \code{gls}, both with
a given workspace and with \code{*workspace_info = 0}, that
\code{*workspace_info} is unchanged after the call, that a workspace
one byte too small gives $\mathit{info} = -5$ and computes nothing, and
that a failing \code{gls} (an indefinite $M$) makes \code{estimate_all}
return $\mathit{info} = 1$ (its call of \code{gls}) at once.

\begin{figure}[t]
\begin{lstlisting}[style=c]
int64_t estimate_all(double *X, double *M, size_t n,
                     int8_t *work, int64_t *workspace_info)
{
    int64_t given = *workspace_info;
    *workspace_info = -1;                         // query the callees
    gls(dummy, dummy, dummy, dummy, work, workspace_info);
    int64_t q1 = *workspace_info;
    *workspace_info = -1;
    postprocess(0, dummy, work, workspace_info);
    int64_t q2 = *workspace_info;
    *workspace_info = given;
    int64_t required = 17600 + up64(q1) + up64(q2);
    if (given == -1) { *workspace_info = required; return 0; }
    if (given > 0 && given < required) return -5;  // info: argument 5
    int8_t *ws = given == 0 ? aligned_alloc(64, required) : work;
    double *y = (double *)(ws + 0), *b = (double *)(ws + 16000);
    int64_t info = 0;
    for (size_t i = 0; i < n; ++i) {
        if (info != 0) continue;   // a failed step: the others do nothing
        observe(i, y);
        *workspace_info = q1;
        if (gls(X, M, y, b, ws + 17600, workspace_info) != 0) { info = 1; continue; }
        *workspace_info = q2;  postprocess(i, b, ws + 17600 + up64(q1), workspace_info);
    }
    *workspace_info = given;
    if (given == 0) free(ws);
    return info;
}
\end{lstlisting}
\caption{The structure of the generated \code{estimate_all}
(\file{examples/compose/compose.mlir}), written out by hand with readable
names; the generated code is in \cref{app:generated}.  \code{dummy} stands
for pointers that a query does not touch, \code{up64} for rounding up to a
multiple of 64.}
\label{lst:compose}
\end{figure}
GLS needs 35.5\,MB (copies of $M$, $X$, $y$ and the matrix $S$); the
driver queries the size once and passes the same workspace to every call.
\file{examples/cases/driver.cc} checks all four cases: the query computes
nothing, a given workspace gives the same result, and a workspace that is
one byte too small gives $\mathit{info} = -5$ and computes nothing; and
that an indefinite matrix assumed spd gives $\mathit{info} = 1$ (its
Cholesky factorisation) and a singular one $\mathit{info} = 1$ (its LU
factorisation), without aborting.

\begin{table}[t]
\centering\small
\begin{tabularx}{\linewidth}{@{}lX@{}}
\toprule
\code{ulm_lamk} op & ulmBLAS call (row-major $m\times n$ storage: strides $(n,1)$) \\
\midrule
\code{gemm} & \code{gemm(m, n, k, alpha, false, A, incRowA, incColA, false, B, ..., beta, C, ...)} \\
\code{gemv} & \code{gemv(m, n, alpha, A, incRowA, incColA, x, incX, beta, y, incY)} \\
\code{trsm} & \code{trlsm} / \code{trusm(m, n, alpha, false, unitDiag, A, ..., B, ...)} \\
\code{trsv} & \code{trlsv} / \code{trusv(n, unitDiag, A, incRowA, incColA, x, incX)} \\
\code{geaxpy} & \code{geaxpy(m, n, alpha, X, incRowX, incColX, Y, incRowY, incColY)} \\
\code{gescal} & \code{gescal(m, n, alpha, X, incRowX, incColX)} \\
\code{potrf} & \code{info = potrf(lower, n, A, incRowA, incColA)} \\
\code{getrf} & \code{info = getrf(m, n, A, incRowA, incColA, piv, incPiv)} (0-based pivots) \\
\code{getrs} & \code{getrs(trans, false, n, nrhs, LU, ..., piv, incPiv, B, ...)} \\
\bottomrule
\end{tabularx}
\caption{The ulmBLAS routines used by \code{convert-ulm-lamk-to-ulmblas}.
\code{trans} selects the strides; \code{trans} and \code{uplo} together
select \code{trl*} or \code{tru*}.}
\label{tab:ulmcalls}
\end{table}

\paragraph{Info.}\label{sec:info}
Numerically, a matrix can fail a property that holds mathematically: a
matrix that is positive definite in exact arithmetic need not be so in
floating point, and the factorisation then breaks down.  LAPACK leaves
the decision to the caller: its routines return \code{info}, and the caller
either chooses another method or returns an \code{info} of its own.
Every function that ulm.lam generates does the same:
\begin{lstlisting}[style=c]
extern "C" int64_t gls(size_t n, size_t m, double *X, double *M, double *y,
                       double *b, int8_t *work, int64_t *workspace_info);
/* returns info:
 *    0  success
 *    1  Cholesky factorisation of v4: not positive definite
 *    2  Cholesky factorisation of a computed matrix: not positive definite
 *   -8  v8: workspace too small
 * (info != 0: the outputs are unspecified) */
\end{lstlisting}
$\mathit{info} = k > 0$ says that check $k$ failed: a factorisation
(\code{potrf}: not positive definite; \code{getrf}: singular) or a call of
another ulm.lam function that returned a non-zero \code{info};
$\mathit{info} = -i$ that argument $i$ is illegal, which is the case for a
workspace that is too small.  The checks are numbered in the order of the
function body, and the interface comment lists them.
\code{ulm-lam-info} implements this on buffers, after
\code{ulm-lam-workspace}: the factorisation kernels return LAPACK's
\code{info}, and after every check the rest of its block runs only if the
check passed (an \code{scf.if} whose result is the status of the block).
A loop carries the status: once an iteration has failed, the remaining
ones do nothing, and the function returns the code of the failure.
Releasing a workspace that the function allocated itself runs also after
a failure.  The guards cost nothing measurable: the Kalman loop of
\cref{sec:kalman} takes the same time as without them.

In Python, a non-zero \code{info} raises \code{lam.InfoError}, whose
message is the line of the interface comment, e.g.
\code{gls: info 1: Cholesky factorisation of v4: not positive definite}.
A fallback, such as an LU factorisation where a Cholesky factorisation
fails, will use the same mechanism as choosing a plan by the sizes at run
time: several plans for one formula in the generated code.

\paragraph{What ulm.lam needs from ulmBLAS next.}
ulm.lam also shows how missing ulmBLAS routines should look: in the style of
the existing ones, with strides for every matrix, a template element type,
\code{bool} flags only where strides cannot express the choice, and
\code{info} as the return value of factorisations.
\begin{itemize}
  \item \emph{The triangle of a product}: \code{sylrk} is used for
        $Z\T Z$ (\cref{sec:storage}); \code{gemmtr} (in ulmBLAS
        \code{master}) would halve the second product of $H P H\T$ and
        $A\T P A$, whose results are symmetric although the factors are
        not transposes of each other.
  \item \emph{QR factorisation} for least squares without normal
        equations, and for full-rank inverses that are not spd:
\begin{lstlisting}[style=c]
template <typename IndexType, typename T>
    IndexType
    geqrf(IndexType m, IndexType n, T *A, IndexType incRowA,
          IndexType incColA, T *tau, IndexType incTau);

//  C := op(Q) C (left) or C op(Q) (right), Q from geqrf
template <typename IndexType, typename T>
    void
    ormqr(bool left, bool trans, IndexType m, IndexType n, IndexType k,
          const T *A, IndexType incRowA, IndexType incColA,
          const T *tau, IndexType incTau,
          T *C, IndexType incRowC, IndexType incColC);
\end{lstlisting}
  \item \emph{Symmetric indefinite factorisation} $PAP\T = LDL\T$
        (Bunch--Kaufman) for inverses of symmetric matrices that are not
        known to be definite, with \code{sytrf} and \code{sytrs} shaped
        like \code{getrf} and \code{getrs} (pivots 0-based, increment
        \code{incPiv}).
\end{itemize}
  The workspace convention of ulm.lam (\cref{sec:workspace}) is the one of
  LAPACK-level routines: BLAS routines never need workspace, LAPACK
  routines such as \code{getri} do.
  A routine that calls other routines with a workspace answers a query
  with the accumulated requirement of its own temporaries and of its
  callees.

Everything the implemented kernel selection emits is already in ulmBLAS:
\code{gemm}, \code{gemv}, \code{trlsm}/\code{trusm},
\code{trlsv}/\code{trusv}, \code{geaxpy}, \code{gescal}, \code{potrf},
\code{getrf} and \code{getrs}.

\paragraph{Reference lowering.}
\code{convert-ulm-lamk-to-linalg} (planned) lowers kernels without a library:
\code{gemm} to \code{linalg.matmul}, triangular solves and factorisations
to \code{scf} loops.
Running it next to \code{convert-ulm-lamk-to-ulmblas} gives differential
tests.

\paragraph{Loops and temporaries.}\label{sec:loops}
Consider matrix operations inside a loop:
\[
  \textbf{for } i = 0,\dots,n-1:\quad B = \mathrm{foo}(i),\quad
  X \mathrel{+}= ABC .
\]
Suppose kernel selection chose $(AB)C$.
The loop body is then a \code{gemm} into a temporary $T=AB$ with
$\beta=0$, followed by a \code{gemm} $X = TC + X$ with $\beta=1$ into the
loop-carried $X$.
\file{examples/loop/} contains this loop, written with \code{ulm_lam} ops as
a frontend would emit it (\cref{lst:loopin}).
\code{ulm-lam-select-kernels} turns it into exactly these two gemms
(\cref{lst:loopsel}).
With upstream passes alone, the result then bufferizes badly:
\begin{itemize}
  \item One-Shot Bufferize allocates $T$ \emph{inside} the loop, and the
        deallocation pass adds run-time ownership bookkeeping to every
        iteration;
  \item $\mathrm{foo}$ allocates $B$ itself, so $B$ is freed in every
        iteration.
        \code{buffer-results-to-out-params} rewrites only functions with a
        body, and \code{foo} is external;
  \item $X$ is returned by value.
        Even with \code{modify-public-functions}, the result becomes a
        separate out parameter, and the updated $X$ is copied into it after
        the loop.
        \code{drop-equivalent-buffer-results} skips public functions.
\end{itemize}

\paragraph{Matrices by reference.}
\code{ulm-lam-results-to-out-params} fixes the calling convention before
bufferization.
It is implemented in \file{lib/UlmLamResultsToOutParams.cpp} and does two
things:
\begin{itemize}
  \item $X = \mathrm{foo}(i)$ is rewritten into $\mathrm{foo}(i, X)$ with
        $X$ passed by reference.
        The external declaration gets one \code{memref} parameter per
        matrix result and returns nothing.
        Every call becomes a \code{memref.alloc}, the call
        \code{foo(i, X)}, and a \code{bufferization.to_tensor restrict
        writable} that makes the written buffer available to the tensor
        program.
        In C, this is the signature \code{void foo(size_t i, double *X)}
        (EmitC maps \code{index} to \code{size_t});
  \item a result marked \code{\{ulm_lam.inout = 2\}} is the updated value
        of argument~2.
        The argument becomes a \code{memref}, it is read once with
        \code{to_tensor}, the final value is written back with
        \code{bufferization.materialize_in_destination}, and the result is
        dropped.
        $X$ is updated in place through the reference.
        Such functions must not be called inside the module; the pass
        rejects that case with a diagnostic.
\end{itemize}
After this pass, the upstream passes \code{one-shot-bufferize},
\code{buffer-loop-hoisting} and \code{buffer-deallocation-pipeline} place
every buffer (\cref{lst:loopin,lst:loop}):
\begin{itemize}
  \item $B$ and $T$ get their buffers once, before the loop: after
        \code{ulm-lam-workspace} they are two slots of the caller's workspace;
  \item $X$ is updated in place;
  \item the loop body contains only the four calls, with no allocation,
        copy or deallocation.
\end{itemize}
\file{examples/loop/check.sh} (\code{make loop}) runs the pipeline and
checks these properties.

\paragraph{The loop in C++.}
The loop example is lowered from \code{ulm_lam} ops to C++ calling ulmBLAS
by passes alone (\file{scripts/ulm-lam-compile.sh}); there is no hand-written
stage.
\Cref{lst:loopc} shows the result: two calls of \code{ulmBLAS::gemm} per
iteration, the second one with $\beta=1$ into $X$, and \code{foo} writing
$B$ into memory that \code{accumulate} allocated and owns.
\file{examples/loop/run.sh} (\code{make loop}) compiles it with
\code{c++ -std=c++17 -Wall -Wextra -pedantic} (no warnings in the generated
code), links it with a C driver that defines \code{foo}, and compares the
result with naive C loops.
For $n=500$ and 4 iterations (one run), it takes 13.2\,ms against
118\,ms for the naive loops; the relative difference is
$1.3\cdot10^{-15}$.
The loop needs a 4\,MB workspace for $B$ and $T$.

\begin{figure}[t]
\lstinputlisting[style=mlir]{\src examples/loop/accumulate.mlir}
\caption{The loop of \cref{sec:loops} as the frontend emits it
(\file{examples/loop/accumulate.mlir}).}
\label{lst:loopin}
\end{figure}

\begin{figure}[t]
\IfFileExists{\src build/loop/accumulate.selected.mlir}{%
  \lstinputlisting[style=mlir]{\src build/loop/accumulate.selected.mlir}}{%
  \fbox{run \code{make loop} to generate this listing}}
\caption{The loop of \cref{lst:loopin} after \code{ulm-lam-normalize},
\code{ulm-lam-infer-properties}, \code{ulm-lam-select-kernels} and
\code{ulm-lam-strip-properties}.
The matrix chain chose $(AB)C$: the first gemm writes $T=AB$ into a new
tensor ($\beta=0$, \texttt{\%cst}), the second accumulates $TC$ into the
loop-carried $X$ (\texttt{\%arg5}, $\beta=1$, \texttt{\%cst\_0}).}
\label{lst:loopsel}
\end{figure}

\begin{figure}[t]
\IfFileExists{\src build/loop/accumulate.bufferized.mlir}{%
  \lstinputlisting[style=mlir,basicstyle=\ttfamily\scriptsize]{\src build/loop/accumulate.bufferized.mlir}}{%
  \fbox{run \code{make loop} to generate this listing}}
\caption{The loop of \cref{lst:loopsel} after
\code{ulm-lam-results-to-out-params}, One-Shot Bufferize,
\code{buffer-loop-hoisting}, deallocation and \code{ulm-lam-workspace}
(generated by \file{scripts/ulm-lam-compile.sh}).
$B$ (\texttt{\%view}) and $T$ (\texttt{\%view\_1}) are slots of the
workspace, set up before the \code{scf.for}; the \code{scf.if} around the
body implements the workspace query and check.
\code{foo} writes $B$ by reference, and $X$ (\texttt{\%arg2}) is updated
in place.}
\label{lst:loop}
\end{figure}

\begin{figure}[t]
\IfFileExists{\src examples/loop/generated/accumulate.cpp}{%
  \srcbetween[style=c]{\src examples/loop/generated/accumulate.cpp}{accumulate(double}{}{loopcpp}}{%
  \fbox{run \code{make loop} to generate this listing}}
\caption{The loop of \cref{lst:loopin} in C++, generated from the
\code{ulm_lam} ops by kernel selection, \code{ulm-lam-results-to-out-params},
upstream bufferization, \code{convert-ulm-lamk-to-ulmblas},
\code{mlir-translate --mlir-to-cpp} and \code{clang-format}: the function
\code{extern "C" int64_t accumulate(...)}, without the \code{\#include}
lines in front of it.
$B$ (\texttt{v18}) and $T$ (\texttt{v20}) live in the workspace; \code{foo}
writes $B$ by reference; $X$ (\texttt{v3}) is updated in place by the
second \code{ulmBLAS::gemm} ($\beta=1$).}
\label{lst:loopc}
\end{figure}

\paragraph{The full pipeline.}
\file{scripts/ulm-lam-compile.sh} runs it in three steps and keeps every
stage as a file:
\begin{lstlisting}[style=sh]
mlir-opt --load-dialect-plugin=UlmLamPasses.dylib \
         --load-pass-plugin=UlmLamPasses.dylib --pass-pipeline='builtin.module(
    func.func(ulm-lam-normalize, ulm-lam-infer-properties, ulm-lam-select-kernels),
    ulm-lam-strip-properties, canonicalize,
    ulm-lam-results-to-out-params, eliminate-empty-tensors,
    one-shot-bufferize{bufferize-function-boundaries
                       function-boundary-type-conversion=identity-layout-map},
    buffer-results-to-out-params{hoist-static-allocs modify-public-functions},
    func.func(buffer-loop-hoisting), buffer-deallocation-pipeline,
    canonicalize, ulm-lam-workspace, canonicalize, cse,
    ulm-lam-info, canonicalize,
    convert-ulm-lamk-to-ulmblas, reconcile-unrealized-casts, form-expressions)'
mlir-translate --mlir-to-cpp gls.emitc.mlir -o gls.cpp
clang-format -i gls.cpp
c++ -std=c++17 -O3 -mcpu=native -isystem $ULMBLAS_DIR -c gls.cpp
\end{lstlisting}
\code{ARCHFLAGS} (default \code{-mcpu=native}) selects the micro kernels of
ulmBLAS, e.g.\ \code{-mcpu=apple-m4} or \code{-mavx -mfma}.

% =====================================================================
\section{Worked example: generalised least squares}\label{sec:example}
% =====================================================================

The example is \cref{eq:gls} with $n=2000$ observations and $m=200$
parameters.
The five files in \file{examples/gls/} show the IR after each step.

\paragraph{Stage 1: frontend output.}
The frontend builds binary products left to right and shares nothing.
Builders infer local properties: $M\inv$ is spd, $X\T$ is full rank.
The binary product $(X\T M\inv)X$ is not yet known to be spd.
\lstinputlisting{\src examples/gls/01-input.mlir}

\paragraph{Stage 2: normal form.}
Products are flattened, and the duplicated $X\T$ and $M\inv$ are shared.
$X\T M\inv X$ is now recognised as spd by the rule $A\T P A$.
\lstinputlisting{\src examples/gls/02-normal.mlir}

\paragraph{Derivation.}
\Cref{tab:derivation} shows Linnea's derivation, which stages 3 and 4
spell out by hand, next to the one \code{ulm-lam-select-kernels} finds from
stage~2 by itself.
Both use Cholesky on $M$, which turns both occurrences of $M\inv$ into
$L\invT L\inv$, compute $W=L\inv X$ once, form $S=W\T W$ from it (the
common subexpression up to transposition), and reuse it for the
right-hand side, $W\T z$ with $z=L\inv y$.
$S$ is spd, so the second Cholesky factorisation turns $S\inv$ into
solves.
The kernel sequences are the same, including \code{syrk} for $W\T W$
(\cref{sec:storage}).

\begin{table}[t]
\centering\small
\begin{tabular}{@{}llr@{\quad}lr@{}}
\toprule
& \multicolumn{2}{c}{Linnea (stages 3, 4)} & \multicolumn{2}{c}{ulm.lam, \code{strategy=gmc}} \\
Step & Kernel & FLOPs & Kernel & FLOPs \\
\midrule
$M=LL\T$ & \code{potrf} & $2.67\cdot10^9$ & \code{potrf} & $2.67\cdot10^9$ \\
$W=L\inv X$ & \code{trsm} & $8.00\cdot10^8$ & \code{trsm} & $8.00\cdot10^8$ \\
$S=W\T W$ & \code{syrk} & $8.00\cdot10^7$ & \code{syrk} & $8.00\cdot10^7$ \\
$S=KK\T$ & \code{potrf} & $2.67\cdot10^6$ & \code{potrf} & $2.67\cdot10^6$ \\
$c=W\T L\inv y$ & \code{trsv}, \code{gemv} & $4.80\cdot10^6$ & \code{trsv}, \code{gemv} & $4.80\cdot10^6$ \\
$b=K\invT K\inv c$ & $2\times$\code{trsv} & $8.00\cdot10^4$ & $2\times$\code{trsv} & $8.00\cdot10^4$ \\
\midrule
total & & $3.55\cdot10^9$ & & $3.55\cdot10^9$ \\
\bottomrule
\end{tabular}
\caption{The derivation for \cref{eq:gls}, with $n=2000$ and $m=200$.
Evaluating \code{inv(X'*inv(M)*X)*X'*inv(M)*y} literally, with
\code{getrf}+\code{getri} ($2n^3$) for each inverse, costs
$1.78\cdot10^{10}$ FLOPs.}
\label{tab:derivation}
\end{table}

\paragraph{Stage 3: kernels on tensors (design).}
\lstinputlisting{\src examples/gls/03-kernels.mlir}

\paragraph{Stage 4: bufferized (design).}
Bufferization inserts exactly three copies, one for each overwritten
input.
$c$, $d$ and $b$ share the result buffer.
\lstinputlisting{\src examples/gls/04-bufferized.mlir}

\paragraph{Stage 5: generated.}
\file{examples/gls/run.sh} compiles stage~2 with
\file{scripts/ulm-lam-compile.sh}.
\Cref{lst:glssel} shows the kernels that \code{ulm-lam-select-kernels}
chooses, \cref{lst:glscpp} the C++ after bufferization and
\code{convert-ulm-lamk-to-ulmblas}.
Bufferization copies exactly $M$, $X$ and $y$, the inputs that kernels
overwrite, and $b$ is written into the caller's buffer.
In the C++, the solve with $K\T$ is \code{ulmBLAS::trusv} on the same
factor $K$ with strides $(1, 200)$, and $W\T$ is the same pointer as $W$
with strides $(1, 200)$, in the \code{gemm} for $W\T W$ and in the
\code{gemv} for $W\T z$.

\begin{figure}[tbp]
\IfFileExists{\src build/gls/gls.selected.mlir}{%
  \lstinputlisting[style=mlir]{\src build/gls/gls.selected.mlir}}{%
  \fbox{run \code{make gls} to generate this listing}}
\caption{GLS after \code{ulm-lam-select-kernels}, generated from stage~2.}
\label{lst:glssel}
\end{figure}

% (not a float: the listing is longer than a page and breaks; the
% environment gives \captionof its box)
\begin{flushleft}
\IfFileExists{\src examples/gls/generated/gls.cpp}{%
  \srcbetween[style=c]{\src examples/gls/generated/gls.cpp}{gls(double}{}{glscpp}}{%
  \fbox{run \code{make gls} to generate this listing}}
\captionsetup{hypcap=false}
\captionof{figure}{GLS in C++ calling ulmBLAS, generated from stage~2: the
function \code{extern "C" int64_t gls(...)}, without the
\code{\#include} lines in front of it.  The outer \code{if}s implement the
workspace convention (\cref{sec:workspace}), the inner ones the
\code{info} of \cref{sec:info}.}
\label{lst:glscpp}
\end{flushleft}

\paragraph{Results.}
\file{driver.cc} uses ulmBLAS for everything:
\begin{itemize}
  \item it builds a random $X$, an spd $M=AA\T+nI$ and a random $y$;
  \item it calls the generated \code{gls};
  \item it verifies the result through the normal equations
        $X\T M\inv(Xb-y)=0$, with $M\inv$ applied by an independent LU
        solve (\code{getrf}, \code{getrs});
  \item it compares against a literal evaluation with explicit inverses
        (\code{getrf}, \code{getri}, \code{gemm}, \code{gemv}).
\end{itemize}
\Cref{tab:results} shows the outcome.
The speedup ($3.7\times$) is smaller than the FLOP ratio ($4.9\times$):
the explicit version consists of large \code{gemm}-like operations that
run at about 140\,GFLOP/s (billions of floating-point operations per
second), while the Cholesky factorisation and the triangular solves of
the derived version run at about 110\,GFLOP/s.

\begin{table}[t]
\centering\small
\begin{tabular}{@{}lr@{}}
\toprule
Kernel sequence chosen by ulm.lam & 29.6\,ms \\
Explicit inverses (same library) & 110.3\,ms \\
Speedup & $3.7\times$ \\
Relative residual $\|X\T M\inv(Xb-y)\|/\|X\T M\inv y\|$ & $1.2\cdot10^{-15}$ \\
$\|b-b_{\text{inv}}\|/\|b\|$ & $2.7\cdot10^{-15}$ \\
\bottomrule
\end{tabular}
\caption{The generated GLS code on an Apple M5 (10 cores), ulmBLAS with
\code{-mcpu=native}, best of 5 runs after warm-up.}
\label{tab:results}
\end{table}

% =====================================================================
\section{Formula (1) in a loop: a Kalman measurement update}\label{sec:kalman}
% =====================================================================

In applications, \cref{eq:gls} is rarely evaluated once.
A sensor array with observation matrix $H$ ($n=2000$ sensors, $m=200$
states) and noise covariance $R$ (spd, correlated) delivers a new
measurement $y_i$ in every time step, and the estimate
\[
  b_i = (H\T R\inv H)\inv H\T R\inv y_i
\]
is needed in every step.
This is the measurement update of a Kalman filter in information form,
$b=(P\inv+H\T R\inv H)\inv(P\inv b_{\text{prior}}+H\T R\inv y)$, with an
uninformative prior ($P\inv=0$).
\file{examples/kalman/estimate.mlir} writes the formula inside the loop,
as on paper (\cref{lst:kalmanin}).

\begin{figure}[t]
\srcbetween[style=mlir]{\src examples/kalman/estimate.mlir}{func.func private @observe}{}{kalmanin}
\caption{Formula (1) in a loop (\file{examples/kalman/estimate.mlir},
without the header comment and the type aliases).
\code{observe} delivers $y_i$, \code{record} consumes $b_i$.}
\label{lst:kalmanin}
\end{figure}

Only $y_i$ changes from step to step, so the pipeline has to do three
things, all of which it does without help:
\begin{itemize}
  \item \code{loop-invariant-code-motion}, an upstream pass, runs before
        kernel selection and hoists the \code{ulm_lam} ops that do not depend
        on $y_i$ out of the loop.
        They are \code{Pure}, so this is safe.
  \item Kernel selection factorises $R$ and $S$ right after their
        definitions, i.e.\ before the loop, and reuses $W=L\inv H$, which
        it computes for $S=W\T W$, for $H\T R\inv y_i = W\T(L\inv y_i)$ in
        the loop.
  \item \code{ulm-lam-results-to-out-params} turns $y_i=\mathrm{observe}(i)$
        into $\mathrm{observe}(i, y)$, and \code{buffer-loop-hoisting}
        moves the buffers of $y$ and $b$ out of the loop.
\end{itemize}
The loop body that results (\cref{lst:kalmanloop}) consists of one
triangular solve with the $2000\times2000$ factor $L$, one \code{gemv}
with $W$ and two triangular solves with the $200\times200$ factor $K$:
$n^2+2nm+2m^2 \approx 4.9\cdot10^6$ FLOPs per step, against
$3.6\cdot10^9$ for evaluating the formula from scratch.
\file{examples/kalman/check.sh} checks this structure.

\begin{figure}[t]
\srcbetween[style=c]{\src examples/kalman/generated/estimate.cpp}{for (size_t}{free(}{kalmanloop}
\caption{The loop body of the generated C++ (\file{scripts/ulm-lam-compile.sh}
on \cref{lst:kalmanin}).  The factorisations, $W$ and $S$ are computed
before the loop.}
\label{lst:kalmanloop}
\end{figure}

\file{examples/kalman/driver.cc} compares five versions
(\cref{tab:kalman}).
Every version is timed with 200 steps and with 0 steps (its one-time part),
best of 3; the table shows the one-time part, the time per step, and the
total for 20 steps.
\begin{itemize}
  \item \emph{ulm.lam}: the pipeline.
  \item \emph{ulm.lam without \code{loop-invariant-code-motion}}: $R$ is still
        factorised once (it is a function argument), but $S$ is formed and
        factorised in every step.
  \item \emph{explicit inverses}: formula (1) evaluated literally in every
        step, with \code{getrf} and \code{getri}.
        This is the same input ulm.lam gets, so the difference to ulm.lam is
        mostly the hoisting that ulm.lam does by itself.
  \item \emph{explicit, hoisted}: written by hand: the invariant inverses
        $R\inv$ and $S\inv$ once before the loop, two \code{gemv} per step.
  \item \emph{operator $G$, hoisted}: written by hand: the operator
        $G=(H\T R\inv H)\inv H\T R\inv$ ($200\times2000$) once, one
        \code{gemv} per step.
\end{itemize}
Every estimate is checked through the normal equations with an
independent LU solve, and all versions agree.

\begin{table}[t]
\centering\small
\begin{tabular}{@{}lrrrr@{}}
\toprule
& before the loop & per step & 20 steps & \\
\midrule
ulm.lam & 15.3\,ms & 0.17\,ms & 18.6\,ms & $1\times$ \\
ulm.lam without \code{loop-invariant-code-motion} & 16.8\,ms & 6.13\,ms & 139.3\,ms & $7.5\times$ \\
explicit inverses in every step & 0.2\,ms & 81.6\,ms & 1631.2\,ms & $88\times$ \\
explicit inverses, hoisted by hand & 69.4\,ms & 0.21\,ms & 73.7\,ms & $4.0\times$ \\
operator $G$, hoisted by hand & 66.5\,ms & 0.24\,ms & 71.2\,ms & $3.8\times$ \\
\midrule
worst normal-equation residual over all steps & \multicolumn{4}{r}{$1.0\cdot10^{-15}$} \\
largest relative difference between the versions & \multicolumn{4}{r}{$2.1\cdot10^{-15}$} \\
\bottomrule
\end{tabular}
\caption{Formula (1) in a loop, $n=2000$, $m=200$, on an Apple M5,
ulmBLAS with \code{-mcpu=native} and OpenMP.  The time per step is
measured over 200 steps (for ulm.lam it includes forming the operator
$P$ before the loop, which pays off from 24 steps on) and includes
\code{observe}, which computes $Hx_i$ plus noise; ``20 steps'' is the
part before the loop plus 20 times the time per step.}
\label{tab:kalman}
\end{table}

\paragraph{The operator before the loop, chosen by the trip count.}
In every step ulm.lam solves with the $2000\times2000$ factor $L$
($O(n^2)$ per step); the hand-written version ``operator $G$'' forms
$G=(H\T R\inv H)\inv H\T R\inv$ once ($O(n^2m)$) and needs one \code{gemv}
($O(nm)$) per step.
Which is better depends on the number of steps.
Kernel selection now decides it: for a product in a loop whose factors
are loop invariant except the last one, it considers the invariant
prefix as an operator $P$ of its own and compares
\[
  \mathrm{cost}(P) + t\cdot 2\,\mathrm{rows}(P)\,\mathrm{cols}(P)\,\mathrm{cols}(y)
  \quad\text{with}\quad t\cdot\mathrm{cost(step)}
\]
for $t$ steps.
$P = S\inv H\T R\inv$ itself is not computable in place (it would solve
on the transposed view $H\T L\invT$ of the reused $W = L\inv H$), but its
transpose is: $P\T = R\inv H S\inv = L\invT W S\inv$, a \code{trsm} on a copy
of $W$ and two \code{trsm} with the factor of $S$; the step is then
$P\T{}\T y$, a \code{gemv} with \code{trans}.
The flops of a step are matrix--vector flops (\code{trsv}, \code{gemv}),
which are memory bound and reach about a tenth of the rate of the
matrix--matrix flops that form $P$ (on the M5 as on Zen~4); they are
weighted with 10.
The trip count is a run-time value here, so the generated code computes
the condition before the loop, forms $P$ if it holds, and runs one of two
loop bodies under the same condition; a constant trip count folds the
branch.
Both variants compute into the same destination, so the loop still
allocates nothing.
For $n=2000$, $m=200$ the operator pays off from 24 steps on, and at 200
steps ulm.lam is faster than both hand-hoisted versions
(\cref{tab:kalman}).

\subsection{The covariance changes in every step}\label{sec:adaptive}

In many applications the noise covariance itself changes: it is
re-estimated online (an adaptive filter), or the noise of the sensors
drifts.
\file{examples/kalman/adaptive.mlir} carries $R$ through the loop and
updates it in every step by calling an external function
(\cref{lst:adaptivein}):
\[
  R_i = \mathrm{update\_covariance}(i, R_{i-1}),\qquad
  b_i = (H\T R_i\inv H)\inv H\T R_i\inv y_i .
\]

\begin{figure}[t]
\srcbetween[style=mlir]{\src examples/kalman/adaptive.mlir}{// update_covariance(i, R) updates}{}{adaptivein}
\caption{The covariance updated in every step
(\file{examples/kalman/adaptive.mlir}, without the header comment and the
type aliases).}
\label{lst:adaptivein}
\end{figure}

\paragraph{In-place updates by reference.}
The result of \code{update_covariance} is marked
\code{\{ulm_lam.inout = 1\}}: the function updates its argument~1 in place.
This extends the by-reference convention of \cref{sec:ulmblas} to
external functions that modify a matrix.
\code{ulm-lam-results-to-out-params} turns the declaration into
\code{void update_covariance(size_t, double *)} and every call into
\code{ulm_lamk.call @update_covariance(\%i) outs(\%R)}.
That op is in destination-passing style, like the kernels, so One-Shot
Bufferize keeps the loop-carried $R$ in one buffer.
A plain call that returns a new matrix would not work here: the loop
would have to yield a new allocation in every step, and deallocating it
needs ownership bookkeeping at run time.
If the matrix is still used after the call, the call updates a copy, so
value semantics are preserved.
\code{ulm_lamk.call} reads and writes memory as far as the compiler knows, so
it is never hoisted, duplicated or removed.

\paragraph{What runs in every step.}
Now $R_i\inv$ changes in every step, so the Cholesky factorisation of
$R_i$, $W=L_i\inv H$, $S=W\T W$ and its factorisation must run in every
step; only what depends on $H$ alone could stay outside.
Kernel selection still shares the factorisation of $R_i$ between both
occurrences of $R_i\inv$ and reuses $W$ for $H\T R_i\inv y_i$
(\cref{lst:adaptiveloop}).
Bufferization copies $R_i$ before the in-place factorisation and $H$
before the in-place solve; the caller's $R_0$
is copied once, before the loop, into the workspace and is not
overwritten.
There is no allocation inside the loop.

The two copies are necessary with in-place kernels.
$R_i$ has two consumers: \code{potrf} overwrites it, and the next call of
\code{update_covariance} reads all of it, so $R_i$ must exist twice.
$H$ is needed in every step, and \code{trsm} overwrites its right-hand
side.
A copy followed by the in-place kernel is what an out-of-place variant
($L := \mathrm{chol}(A)$, $W := L\inv H$) would do; ulm.lam generates it
inline where an input is still needed, so ulmBLAS keeps only the in-place
kernels.
With the factorisation at about 25\,ms per step, the copies are a small
part of the time.

\begin{figure}[t]
\srcbetween[style=c]{\src examples/kalman/generated/adaptive.cpp}{for (size_t}{free(}{adaptiveloop}
\caption{The loop body generated from \cref{lst:adaptivein}: an in-place
update of $R$, two copies, two Cholesky factorisations, one
\code{trlsm}, and the solves of \cref{lst:kalmanloop}.}
\label{lst:adaptiveloop}
\end{figure}

\file{examples/kalman/driver_adaptive.cc} defines an
\code{update_covariance} that moves $R$ towards a base covariance and
changes the noise level with the step, which keeps $R$ spd.
Over 10 steps, ulm.lam's version takes 338\,ms (34\,ms per step, dominated by
the Cholesky factorisation of the $2000\times2000$ matrix $R_i$),
against 1177\,ms for formula (1) with explicit inverses on the same
covariances ($3.5\times$).
Its workspace is 67.5\,MB: the copy of $R$ that is updated in place,
the factor $L_i$, $W$, $S$, $y_i$ and $b_i$.
The worst normal-equation residual over all steps is
$9.8\cdot10^{-16}$, and both versions agree to $2.1\cdot10^{-15}$.

A full Kalman filter also updates the state covariance $P$ and predicts
with the state transition.
In information form the measurement update is a sum,
$P\inv \leftarrow P\inv + H\T R\inv H$, so it needs inverses as terms of
sums and the identity matrix, which kernel selection does not handle yet.

% =====================================================================
\section{Python frontend}\label{sec:python}
% =====================================================================

\code{python/ulm/lam/} lets the linear algebra be written in Python and
compiled by ulm.lam (\cref{lst:pykalman}): \code{from ulm import lam}.
It has two parts: a frontend that turns Python into \code{ulm_lam} IR, and
a runtime that compiles it and calls the result.

\begin{figure}[t]
\srcbetween[style=py]{\src python/examples/kalman.py}{observe = lam.extern}{def best(}{pykalman}
\caption{Formula (1) in a loop, in Python (from
\file{python/examples/kalman.py}).  \code{observe} and \code{record} are
external functions, implemented in Python by callbacks.}
\label{lst:pykalman}
\end{figure}

\paragraph{The language.}
A function decorated with \code{@lam.function} is traced: it is called with
symbolic values, and the Python operators build \code{ulm_lam} operations
(\code{@} is \code{ulm_lam.mul}, \code{+} and \code{-} are \code{ulm_lam.add},
a number times a matrix is \code{ulm_lam.scale}, \code{.T} is
\code{ulm_lam.transpose}, \code{.inv} is \code{ulm_lam.inverse}).
\begin{itemize}
  \item Parameters have types with properties: \code{Matrix(n, n, "spd")},
        \code{Vector(n)}, \code{Index}; a symmetric matrix can be stored
        as one triangle, \code{Matrix(n, n, "spd", stored="lower")}
        (\cref{sec:storage}).
  \item Sizes are numbers or names: \code{Matrix("n", "m", "full_rank")}
        has run-time sizes (\cref{sec:dynamic}).  Equal names are equal
        sizes; they become index parameters in front of the matrices, in
        order of first appearance.  From Python they are taken from the
        arrays (and checked: \code{gls: M: size n = 4, but 5 from an
        earlier argument}); external functions get the sizes they use as
        parameters too, and Python callbacks receive arrays of the actual
        shapes.  One compiled library serves all sizes.
  \item The workspace: every library has one (\code{lib.workspace},
        64-byte aligned), which its functions reuse and which grows as
        needed; \code{lib.f(..., work=lam.Workspace(size))} passes another
        one (\code{lib.f.workspace_size(...)} is the query).  A call from
        a Python callback gets its own.
        Intermediate results get their properties from
        \code{ulm-lam-infer-properties}.
  \item Parameters have intents: a matrix is \code{in} by default;
        \code{InOut(Matrix(...))} may be overwritten, and with
        \code{lam.update(param, value)} it holds \code{value} on exit.
        Return values are \code{out}.
  \item \code{Float} parameters are scalars known only at run time; they
        scale matrices (\code{s2 * A}).  \code{lam.eye(n)} is the
        identity ($n$ a number or a size name), never formed: a term
        $\alpha I$ of a sum is added to the diagonal.
        \code{lam.assume(x, "spd")} asserts properties the compiler cannot
        derive (e.g.\ for $K + \sigma^2 I$); a factorisation checks them at
        run time (\code{info}).
        So far a run-time scalar may scale identities only; a scaled
        matrix or product with one is the next step.
  \item \code{with lam.loop(n, carry=[X]) as L:} is a loop with
        loop-carried matrices (\code{L.carried}, \code{L.next(...)},
        \code{L.results}); the formula can be written inside the loop, and
        \code{loop-invariant-code-motion} and kernel selection move what does
        not depend on the loop out of it.
  \item \code{lam.extern(name, params, result, updates=k,
        workspace=True)} declares an external function; a lamp function
        can also call another one, and their workspaces accumulate.
  \item Shapes are checked while tracing: a dimension mismatch, or a loop
        value used outside its loop, is a Python exception at the line that
        causes it.
\end{itemize}
The frontend writes the IR as text in the syntax of the dialects; it needs
neither the MLIR Python bindings nor any package beyond the Python standard
library.

\paragraph{The runtime.}
\code{lam.compile(functions)} writes the module, runs
\file{scripts/ulm-lam-compile.sh}, and compiles the C++ with ulmBLAS into a
shared library, which it loads with \code{ctypes}.
An unchanged program reuses its library, so recompiling is only needed
after a change.
Each function of the library is a Python callable whose signature comes
from the interface comment the compiler generates
(\cref{sec:workspace}):
\begin{itemize}
  \item \code{in} matrices are any C-contiguous \code{float64} buffer of
        the right size (NumPy arrays, \code{array.array}, nested lists);
        other layouts are copied;
  \item \code{inout} matrices must be writable and C-contiguous, because
        the function may overwrite them;
  \item \code{out} matrices are allocated by the runtime and returned
        (NumPy arrays if NumPy is installed);
  \item the workspace is queried (\code{workspace_info = -1}) and
        allocated by the runtime.
\end{itemize}
The generated C++ is available from Python: \code{lib.estimate.cpp} is
the code of one function, with its interface comment, and
\code{lib.source} the whole file; the examples print it.
External functions can be implemented in Python
(\code{lib.bind(observe=..., record=...)}).
The compiler generates a C trampoline for each of them; matrix arguments
reach the callback as NumPy views of the caller's memory (or memoryviews),
a returned matrix is copied into the caller's buffer, and an exception in a
callback is raised again when the compiled function returns.

\paragraph{Results.}
\file{python/examples/kalman.py} runs the loop of \cref{lst:pykalman}
with $n=4000$ sensors, $m=200$ states and 200 steps, measurements and
estimates passing through Python callbacks, and compares it with NumPy
(linked to OpenBLAS):

\begin{center}
\small
\begin{tabular}{@{}lr@{}}
\toprule
ulm.lam (generated C++, ulmBLAS) & 122.0\,ms \\
NumPy, the formula as written in the loop & 270271.2\,ms \\
NumPy, the operator $G$ computed by hand before the loop & 368.4\,ms \\
\midrule
relative difference to NumPy & $4.1\cdot10^{-15}$ \\
\bottomrule
\end{tabular}
\end{center}

\noindent
Written as on paper, the Python program is $3\times$ faster than the
version a NumPy programmer optimises by hand, and more than 2000 times
faster than the formula written literally.
Before the operator was formed by kernel selection (with the per-step
solves), ulm.lam took 482\,ms here and lost against the hand-written
version.
\file{python/tests/test_lam.py} (44 tests, \code{make python}) checks
GLS with lists, buffers and NumPy arrays, \code{inout} parameters, loops
with carried matrices and callbacks, the operator chosen by the trip
count, functions calling functions, two results, NumPy layouts, one
triangle of symmetric matrices, run-time sizes, identity terms with
run-time scalars, \code{info},
exceptions in callbacks, and the errors of the frontend.

\paragraph{The adaptive filter in Python.}
\file{python/examples/adaptive_kalman.py} is the adaptive filter of
\cref{sec:adaptive}: the covariance $R$ is carried through the loop and
updated in place in every step by \code{update_covariance}, declared
with \code{updates=1} and implemented in Python, where it changes $R$
through a NumPy view of the generated code's memory
(\cref{lst:pyadaptive}).

\begin{figure}[t]
\srcbetween[style=py]{\src python/examples/adaptive_kalman.py}{update_covariance = lam.extern}{def main(}{pyadaptive}
\caption{The adaptive filter in Python (from
\file{python/examples/adaptive_kalman.py}): $R$ is loop-carried and updated
in place by an external function, a Python callback.}
\label{lst:pyadaptive}
\end{figure}

With $n=4000$, $m=200$ and 10 steps the Cholesky factorisation of $R_i$
dominates, and nothing expensive can move out of the loop.
NumPy has no Cholesky solve, so the best one can write by hand is one LU
solve per step for $[H, y_i]$:

\begin{center}
\small
\begin{tabular}{@{}lr@{}}
\toprule
ulm.lam (generated C++, ulmBLAS) & 1758.7\,ms \\
NumPy, the formula as written in the loop & 8099.1\,ms \\
NumPy, one LU solve per step, by hand & 2888.8\,ms \\
\midrule
relative difference to NumPy & $7.9\cdot10^{-15}$ \\
\bottomrule
\end{tabular}
\end{center}

\noindent
ulm.lam is $1.6\times$ faster than the hand-written NumPy version
(Cholesky instead of LU, and $W=L_i\inv H$ reused for $H\T R_i\inv y_i$)
and $4.6\times$ faster than the formula as written.
All three times include the Python callback that updates $R$ ($O(n^2)$
per step), and the caller's $R_0$ (an \code{in} parameter) is not
overwritten.

% =====================================================================
\section{Julia frontend}\label{sec:julia}
% =====================================================================
\file{julia/UlmLam} is a Julia package with the same frontend as the
Python one: functions are traced to \code{ulm_lam} IR, compiled by the
same pipeline, and called with \code{ccall}; no Python is involved.
The notation is Julia's:
\begin{lstlisting}[style=py]
using UlmLam, LinearAlgebra

@lam function gls(X::Mat(:n, :m, :full_rank), M::Mat(:n, :n, :spd), y::Vec(:n))
    inv(X' * inv(M) * X) * X' * inv(M) * y
end

@lam function gp(K::Mat(:n, :n, :symmetric), s2::Float, Ks::Mat(:n, :m),
                 Kss::Mat(:m, :m, :symmetric), y::Vec(:n))
    Kn = assume(K + s2 * I, :spd)
    (Ks' * inv(Kn) * y, Kss - Ks' * inv(Kn) * Ks)
end

b = gls(X, M, y)                # Julia arrays; compiled at the first call
@code_lam gls(X, M, y)          # the derivation
\end{lstlisting}
\code{X'}, \code{inv}, \code{A \textbackslash{} B}, \code{B / A},
\code{I} from \code{LinearAlgebra}, \code{assume}, \code{loop} (a
do-block, with loop-carried matrices as its further arguments and
results), \code{extern} with callbacks bound by \code{bind!}, \code{InOut},
\code{Float} and \code{Index} parameters, sizes as symbols, the workspace,
\code{InfoError} and both backends work as in Python.

\paragraph{Column-major arrays without copies.}
Julia stores matrices by columns, the generated code by rows.
A Julia matrix $X$ ($n \times m$) is, as memory, the row-major matrix
$X\T$ ($m \times n$).
So the frontend declares a matrix parameter with the transposed type
(triangles flipped: a lower triangular $X$ is an upper triangular $X\T$),
the traced function sees \code{transpose(P)}, and a matrix result is
returned transposed; vectors are the same in both orders.
Kernel selection turns these transposes into strides, so no matrix is
copied for the conversion: GLS from Julia gets exactly the kernel
sequence of GLS from Python.
That needed four general improvements of the selector, which help
Python programs as well:
transposes are pushed to the leaves in the normal form
($(AB)\T = B\T A\T$, also for sums and scales), and an \code{assume} that
only drops properties no longer hides a value from the rules (so $P$ and
$P\T$ are recognised as transposes of each other for \code{syrk} and
symmetry);
a solve from the right with an LU factorisation is \code{getrs} on the
transposed view of the right-hand side (in place);
an operand of an in-place solve that is a transposed view is copied
transposed (\code{gecopy}, $O(nm)$, only if nothing cheaper exists);
and a transpose used outside the algebra (a loop's initial value, a
result) is materialised by \code{gecopy}.
Callbacks of external functions get vectors as \code{Vector} and matrices
as lazy transposes of the row-major buffers, without copies.

\paragraph{Calling, inspecting, Julia's BLAS.}
A \code{@lam} function called with Julia arrays is compiled at its first
call (about half a second: the pipeline takes 0.3\,s, the C++ compiler
the rest) and the library is kept; \code{UlmLam.compile} builds one
library for several functions, \code{bind!(f; \ldots)} implements the
external functions of \code{f}.
\code{@code_lam f(args)} shows the derivation: the kernels after selection
(\cref{sec:select}), each written as an assignment in place in Julia's
terms (\code{T1 ← L1⁻¹ T1} for a \code{trsm}), with loops, the plans of a
decision taken at run time, callbacks, and the meaning of every
\code{info} value.
Since the IR holds a Julia matrix as its transpose, each IR value is
shown as a matrix or its transpose; where the IR leaves the choice (a new
temporary), \code{@code_lam} shows the orientation with the fewest
transposes, so that $Q (K\T V)$ reads as \code{T1 ← Kᵀ V},
\code{T2 ← Q T1}.
The default backend, \code{:julia}, is the CBLAS lowering compiled
against the libblastrampoline Julia has loaded, with 64-bit integers and
its \code{64\_} symbols (\file{include/ulm\_lam/cblas\_backend.h},
\code{ULM\_LAM\_CBLAS\_ILP64}): the generated code runs on OpenBLAS, or on
ulmBLAS after \code{UlmBLAS.forward()} of
UlmBLAS.jl~\cite{ulmblasjl}; \code{:ulmblas} and \code{:cblas} are the
other two.
The Pluto notebook \file{notebook/ulm-lam.jl} demonstrates all of this
(GLS, GP regression against AbstractGPs.jl, the Kalman filter, linear
attention).

\paragraph{Vector spaces.}
\file{julia/UlmLam/test/vectorspace.jl} checks 137 expressions on
$\mathbb{R}^n$ and $\mathbb{R}^{m\times n}$ against Julia: sums,
differences, negation, constant and run-time scalars on either side,
division by scalars and linear combinations, combined with products,
transposes, solves and inverses.
Making all of them work needed four additions to kernel selection: a
constant scale of a value used outside the algebra (a result, \code{x *
2} after a solve) is a \code{gescal}; a transposed term of a sum is copied
transposed before the \code{geaxpy}s; an inverse used as a value (a
result, a term of a sum, also scaled) is formed as $A^{-1} = A^{-1} E$
with $E$ the identity built by \code{gescal} with 0 and \code{adddiag},
that is, by the factorisation and a solve that any $A^{-1}B$ gets (in
products it is still never formed); and \code{gescal} with $\alpha = 0$
sets the matrix to zero, as BLAS's \code{scal} does (also NaN).
The frontend divides by run-time scalars ($x / a$, $a \backslash x$).

\paragraph{Results.}
\file{julia/UlmLam/test/runtests.jl} (75 tests, and the 137 above) runs GLS with fixed and
run-time sizes, solves from the left and from the right (Cholesky and
LU), a chain chosen at run time, two results, a loop with callbacks, a
loop-carried matrix, GP with \code{s2 * I} and a wrong assumption
(\code{InfoError}), \code{InOut}, a given workspace and argument errors,
on Julia's BLAS, on ulmBLAS and on OpenBLAS (\code{backend=:cblas}),
and \code{@code_lam}.
\file{julia/examples/gp.jl} and \file{julia/examples/kalman.jl} are the
benchmarks of \cref{sec:glance}: the GP prediction against Julia code
written by hand ($2$--$2.4\times$ faster on ulmBLAS, $1.3$--$1.6\times$ on
the same OpenBLAS) and the Kalman loop ($3$--$4\times$ faster than the
operator formed by hand).

% =====================================================================
\section{Showcase: Gaussian process regression}\label{sec:gp}
% =====================================================================
A GP (Gaussian process) with kernel $k$ predicts at test points $X_*$ from
training data $(X, y)$ by
\[
  \mu = K_*\T K\inv y, \qquad \Sigma = K_{**} - K_*\T K\inv K_*,
\]
with $K = k(X, X) + \sigma^2 I$ ($n\times n$, spd), $K_* = k(X, X_*)$
($n\times m$) and $K_{**} = k(X_*, X_*)$.
Written literally in NumPy, this forms $K\inv$.
Libraries therefore implement it by hand; scikit-learn's
\code{GaussianProcessRegressor}, for instance, computes
$L = \mathrm{chol}(K)$, $\alpha = L\invT L\inv y$, $V = L\inv K_*$,
$\mu = K_*\T\alpha$ and $\Sigma = K_{**} - V\T V$ in \code{fit} and
\code{predict}.
That is the hand work of FLENS-LAPACK again: someone had to notice that
$K_*\T K\inv K_* = (L\inv K_*)\T (L\inv K_*)$.
In ulm.lam the formula is the program
(\code{python/examples/gp.py}):
\begin{lstlisting}[style=py]
@lam.function
def gp(K: Matrix("n", "n", "symmetric"), s2: Float, Ks: Matrix("n", "m"),
       Kss: Matrix("m", "m", "symmetric"), y: Vector("n")):
    Kn = lam.assume(K + s2 * lam.eye("n"), "spd")
    mean = Ks.T @ Kn.inv @ y
    cov = Kss - Ks.T @ Kn.inv @ Ks
    return mean, cov
\end{lstlisting}
Here $K$ is the kernel matrix $k(X, X)$ and \code{s2} the noise variance
$\sigma^2$, a \code{Float} parameter known only at run time.
\code{lam.eye("n")} is the identity of the run-time size $n$; it is never
formed: a term $\alpha I$ of a sum becomes the kernel \code{adddiag}
($A := A + \alpha I$, $O(n)$), with $\alpha$ a constant or a run-time
value.
Whether $K + \sigma^2 I$ is spd depends on the sign of $\sigma^2$, which
the compiler cannot know, so the program asserts it with
\code{lam.assume}; the Cholesky factorisation checks the assertion at run
time (a negative \code{s2} gives \code{info} 1, ``Cholesky factorisation
of a matrix computed from v3: not positive definite'').
Kernel selection derives the hand version: \code{adddiag} on the copy of
$K$ that is factorised, one \code{potrf}, two
\code{trsv} and a \code{gemv} for $\mu$, one \code{trlsm} for
$V = L\inv K_*$, and \code{sylrk} with $\beta = 1$ into a copy of
$K_{**}$ for $\Sigma$, which is then completed (\code{sycomplete}).
For the last step the \code{syrk} rule of \cref{sec:storage} was extended:
a Gram product $Z\T Z$ may also be added to a symmetric matrix whose lower
triangle is valid ($C = \beta C + \alpha Z\T Z$).
The sizes are run-time names, so one library serves all $n$ and $m$.
\Cref{lst:gpcpp} in \cref{sec:glance} shows the generated function.
With \code{backend="cblas"} (\cref{sec:cblas}) the code is the same, with
\code{ulm_lam_cblas::} instead of \code{ulmBLAS::}
(\file{python/examples/generated/gp\_cblas.cpp}).



\paragraph{Data.}
Two public data sets, downloaded on first use:
the weekly CO$_2$ concentration at Mauna Loa (NOAA, National Oceanic and
Atmospheric Administration, and Scripps; 2718 weeks since 1974), with
the kernel and the hyperparameters of
Rasmussen and Williams (Gaussian Processes for Machine Learning, 2006,
Sec.~5.4.3), and the hourly rentals of a bike sharing system (UCI,
University of California, Irvine, Machine Learning Repository; 17\,379
hours, 14 features).
For CO$_2$, a back test trained until 2016 predicts the following ten
years with an RMSE (root mean square error) of 2.8\,ppm, and a forecast
from all data reaches 2036 (\cref{fig:co2}; the kernel of 2006
underestimates the accelerating trend).
The results agree with SciPy and scikit-learn to $2.5\cdot 10^{-11}$.

\IfFileExists{../build/gp/co2.png}{%
\begin{figure}[t]
\centering
\includegraphics[width=\linewidth]{../build/gp/co2.png}
\caption{CO$_2$ at Mauna Loa: back test and forecast by \code{gp}
(mean $\pm 2$ standard deviations).}\label{fig:co2}
\end{figure}}{}

\paragraph{Timings.}
Bike sharing, $m = 1000$ test hours, Apple M5, without computing the kernel
matrices (scikit-learn: with them, since it computes them inside).
ulm.lam runs twice: on ulmBLAS (with OpenMP, so \code{potrf} and
\code{trtrs} run as tasks), and with \code{backend="cblas"}
(\cref{sec:cblas}) on OpenBLAS, the BLAS of NumPy and SciPy here, so
that the comparison with the hand-written versions measures the mapping
alone:

\begin{center}
\begin{tabular}{rrrrrr}
\toprule
$n$ & ulm.lam/OpenBLAS & ulm.lam/ulmBLAS & SciPy by hand & NumPy literal & scikit-learn \\
\midrule
1000 & 0.018\,s & 0.015\,s & 0.018\,s & 0.052\,s & 0.048\,s \\
2000 & 0.052\,s & 0.045\,s & 0.068\,s & 0.287\,s & 0.160\,s \\
4000 & 0.240\,s & 0.182\,s & 0.313\,s & 1.589\,s & 0.654\,s \\
8000 & 1.619\,s & 1.189\,s & 1.610\,s & 9.929\,s & 3.531\,s \\
\bottomrule
\end{tabular}
\end{center}

On the same BLAS, ulm.lam is as fast as SciPy by hand or up to
$1.3\times$ faster, $1.5$--$1.8\times$ faster than scikit-learn and
$3$--$7\times$ faster than the literal formula, and it gets the hand
version without anyone writing it.
On ulmBLAS it is another $1.3$--$1.4\times$ faster ($n \ge 4000$): on the
M5, ulmBLAS's \code{potrf} with OpenMP beats OpenBLAS (at $n = 8000$:
0.54\,s against 0.67\,s).

\paragraph{What the example shows is missing.}
The kernel matrices $k(X, X)$, $k(X, X_*)$ and $k(X_*, X_*)$ are
parameters of \code{gp}: the Python script \file{gp.py} computes them
before the call, with NumPy (elementwise functions of the data, e.g.\
$\exp(-\lVert x_i - x_j\rVert^2/2\ell^2)$).
To move them into the ulm.lam program, ulm.lam would need elementwise
operations; it could then compute only the triangle that \code{potrf}
reads, obtain squared distances from \code{syrk}, and write $K$ directly
into the buffer that is factorised.
Likewise the hyperparameters ($\ell$, $\sigma$, \dots) are fitted
outside, by scikit-learn on 1500 points: fitting needs
$\log\det K = 2\sum_i \log L_{ii}$ and traces such as
$\operatorname{tr}(K\inv \partial K)$, scalar-valued operations that would
reuse the Cholesky factor $L$ that \code{gp} already computes.
Often only the variances $\operatorname{diag}(\Sigma)$ are needed; after
$V = L\inv K_*$ they cost $O(nm)$ instead of the $O(nm^2)$ of the whole
$\Sigma$.

\subsection{Large problems: $n = 100\,000$ on bwUniCluster 3.0}\label{sec:gpcluster}
\file{examples/gp\_cluster} runs the same function with $K$ as
\code{inout} (factorised in place: at $n = 100\,000$, $K$ has 80\,GB) on the
3D Road Network data (UCI; 434\,874 points, the altitude predicted from the
position, $m = 2000$, Mat\'ern 3/2 kernel), compiled on the cluster for its
Zen~4 nodes (2$\times$ AMD EPYC 9454, 96 cores; GCC 14.2,
\code{-march=native}, OpenMP).
It compares ulm.lam on ulmBLAS, ulm.lam on MKL (\code{backend=cblas}) and
the same steps written by hand with MKL, as SciPy on MKL would do them
(times in seconds, 96 threads):

\begin{center}
\begin{tabular}{rrrrr}
\toprule
$n$ & RMSE & ulm.lam/ulmBLAS & ulm.lam/MKL & by hand/MKL \\
\midrule
4\,000 & 11.03\,m & 0.44 & 0.17 & 0.16 \\
12\,500 & 6.85\,m & 2.48 & 1.34 & 1.46 \\
25\,000 & 4.83\,m & 13.66 & 5.71 & 5.85 \\
50\,000 & 3.17\,m & 106.5 & 26.15 & 27.84 \\
100\,000 & 1.95\,m & 590.7 & 153.7 & 150.7 \\
\bottomrule
\end{tabular}
\end{center}

\emph{The mapping} is as fast as the hand-written version on the same BLAS
(within 6\,\% either way).
In the first measurements, whichever of the two ran first in a process
paid about 25\,\%: the driver zeroed its arrays with the main thread
(\code{std::vector}), so all pages lay on one of the eight NUMA domains,
and the kernel's automatic NUMA balancing migrated them during the first
heavy run.
With the arrays first written by the threads that compute them (the table
above), the order no longer matters ($n = 50\,000$, both orders: ulm.lam
26.0 and 26.4\,s, by hand 27.4\,s; \code{numactl --interleave=all} has the
same effect), and the kernel matrices take 1.1\,s instead of 26\,s at
$n = 100\,000$.
\emph{ulmBLAS} is $3.8\times$ slower than MKL at $n = 100\,000$, although
it is faster single threaded ($n = 12\,500$: 18.3\,s against 22.6\,s):
its task-parallel \code{potrf} and \code{trtrs} scale to about 12 threads
(one NUMA domain) and hardly beyond, and \code{sylrk} runs on one thread.
The error of the prediction falls from 15.5\,m ($n = 1000$) to 1.95\,m:
exact GP regression profits from all the data, which makes the
$O(n^3)$ factorisation, and a fast \code{potrf}, the point.

% =====================================================================
\addtocontents{toc}{\protect\setcounter{tocdepth}{2}}
\section{More showcases}\label{sec:showcases}
% =====================================================================
Five more programs, each written as the formula and compared with NumPy
(and one with PyTorch)
(OpenBLAS) on the Apple M5 (\code{make showcases}).

\subsection{Ridge regression}\label{sec:ridge}
\file{python/examples/ridge.py}:
\begin{lstlisting}[style=py]
@lam.function
def ridge(X: Matrix("n", "p"), Y: Matrix("n", "k"), lam_: Float):
    A = lam.assume(X.T @ X + lam_ * lam.eye("p"), "spd")
    return A.inv @ X.T @ Y
\end{lstlisting}
on the bike-sharing data, $n = 16\,000$ training hours, with $p$ random
Fourier features (Rahimi and Recht: ridge regression on them approximates
kernel ridge regression with an RBF kernel).
ulm.lam computes $X\T X$ by \code{syrk} (one triangle), adds $\lambda$ to
its diagonal, factorises the triangle by Cholesky, and chooses the order
of $A\inv X\T Y$ at run time from $n$, $p$ and $k$.
Times in seconds:

\begin{center}
\begin{tabular}{rrrrrrr}
\toprule
$p$ & RMSE & ulm.lam & NumPy with \code{inv} & NumPy \code{solve} & SciPy Cholesky & scikit-learn \\
\midrule
500 & 0.413 & 0.023 & 0.066 & 0.022 & 0.023 & 0.026 \\
1000 & 0.327 & 0.075 & 0.318 & 0.075 & 0.074 & 0.084 \\
2000 & 0.317 & 0.280 & 0.862 & 0.300 & 0.282 & 0.336 \\
4000 & 0.292 & 1.290 & 3.679 & 1.422 & 1.341 & 1.314 \\
\bottomrule
\end{tabular}
\end{center}

\noindent
The time is that of $X\T X$ ($np^2$), which NumPy also computes by
\code{syrk}: written as the formula, ulm.lam is as fast as the best
version written by hand, and $3$--$4\times$ faster than the formula with
\code{inv}.

\subsection{Linear attention}\label{sec:attention}
\file{python/examples/attention.py}:
\begin{lstlisting}[style=py]
@lam.function
def attention(Q: Matrix("N", "d"), K: Matrix("N", "d"), V: Matrix("N", "e")):
    return Q @ K.T @ V
\end{lstlisting}
Linear attention (Katharopoulos et al., 2020) replaces
$\mathrm{softmax}(QK\T)V$ by $\phi(Q)\phi(K)\T V$, normalised by
$\phi(Q)\phi(K)\T\mathbf{1}$ (a column of ones appended to $V$; $\phi$ and
the division are elementwise and done in NumPy).
Its point is the order: $\phi(Q)(\phi(K)\T V)$ costs $O(Nd^2)$, the order
in which NumPy evaluates the formula $O(N^2 d)$.
ulm.lam chooses it at run time from $N$ and $d$ (head dimension
$d = 64$, seconds):

\begin{center}
\begin{tabular}{rrrr}
\toprule
$N$ & ulm.lam & NumPy as written & NumPy, parentheses by hand \\
\midrule
256 & 0.0001 & 0.0002 & 0.0001 \\
1024 & 0.0001 & 0.0017 & 0.0003 \\
4096 & 0.0005 & 0.0222 & 0.0010 \\
16384 & 0.0031 & 0.5053 & 0.0026 \\
\bottomrule
\end{tabular}
\end{center}

\noindent
Up to $160\times$ faster than the formula in NumPy, as fast as the
parentheses placed by hand; and for a short sequence with a large
dimension ($N = 16$, $d = 512$) ulm.lam takes the other order, which is
$8\times$ faster there than the one that is right for long sequences.

\subsection{The Kalman filter}\label{sec:kalmanfilter}
\file{python/examples/kalman_filter.py}: predict and update in a loop, as
in the textbook,
\begin{lstlisting}[style=py]
with lam.loop(steps, carry=[x0, P0]) as L:
    x, P = L.carried
    P = lam.assume(P, "symmetric")
    x = F @ x                                       # predict
    P = F @ P @ F.T + Q
    S = lam.assume(H @ P @ H.T + R, "spd")          # update
    K = P @ H.T @ S.inv
    x = x + K @ (measure(L.index) - H @ x)
    P = P - K @ H @ P
    estimate(L.index, x)
    L.next(x, P)
\end{lstlisting}
for data assimilation: a field on a periodic grid of $n$ points,
transported and diffused, observed by $m$ sensors.
Per step ulm.lam generates the expert's version: $Fx$ (\code{gemv}),
$FPF\T + Q$ (two \code{gemm}), $HP$ (\code{gemm}) and
$S = (HP)H\T + R$, the Cholesky factorisation of $S$,
$K = (HP)\T S\inv$ by two triangular solves on a transposed copy of the
computed $HP$ ($P$ is symmetric, so $PH\T = (HP)\T$), the update of $x$
(two \code{gemv}), and $P - K(HP)$ (one \code{gemm}, $n^2 m$).
NumPy evaluates $K H P$ from left to right, $(KH)P$, which adds $n^3$;
the expert writes $HP$ once and solves with Cholesky (seconds):

\begin{center}
\begin{tabular}{rrrrrr}
\toprule
$n$ & $m$ & steps & ulm.lam & NumPy as written & NumPy by hand \\
\midrule
500 & 50 & 50 & 0.139 & 0.348 & 0.155 \\
1000 & 100 & 50 & 0.927 & 1.673 & 1.064 \\
2000 & 200 & 20 & 2.623 & 4.550 & 3.094 \\
\bottomrule
\end{tabular}
\end{center}

\subsection{Bayesian linear regression}\label{sec:bayes}
\file{python/examples/bayes.py}: with the prior $w \sim N(0, \alpha\inv I)$ and the noise precision
$\beta$ (Bishop, Sec.~3.3),
\begin{lstlisting}[style=py]
@lam.function
def blr(Phi: Matrix("n", "p"), t: Vector("n"), Phis: Matrix("m", "p"),
        alpha: Float, beta: Float):
    A = lam.assume(alpha * lam.eye("p") + beta * (Phi.T @ Phi), "spd")
    w = beta * A.inv @ Phi.T @ t
    return Phis @ w, Phis @ A.inv @ Phis.T + (1 / beta) * lam.eye("m")
\end{lstlisting}
$\alpha$ and $\beta$ are scalars known at run time, which scale products
here and enter arithmetic ($1/\beta$).
Kernel selection collects such factors of a product, computes the product
without them and scales the result (\code{gescal}); scaling commutes with
the product, so $\beta$ is applied to the predictive mean ($m$ values),
not to the weights.
A run-time scale of a matrix in a sum is a \code{gescal}, of the identity
an \code{adddiag}; a symmetric matrix stored as one triangle stays so.
The generated sequence: $\Phi\T\Phi$ by \code{syrk}, scaled by $\beta$,
$\alpha$ added to the diagonal, Cholesky of the triangle; the mean by
\code{gemv}, two \code{trsv} and \code{gemv}; the covariance as $V\T V$
with $V = L\inv\Phi_*\T$ (\code{trsm}, \code{syrk}) and $1/\beta$ added to
its diagonal.
On the bike-sharing data with random Fourier features, the predictive
mean and the full covariance of 1379 test hours (seconds):

\begin{center}
\begin{tabular}{rrrrrr}
\toprule
$p$ & RMSE & in 2 sd & ulm.lam & NumPy as written & SciPy by hand \\
\midrule
500 & 0.404 & 89\,\% & 0.030 & 0.138 & 0.028 \\
1000 & 0.314 & 95\,\% & 0.093 & 0.334 & 0.089 \\
2000 & 0.286 & 97\,\% & 0.336 & 1.135 & 0.333 \\
4000 & 0.259 & 98\,\% & 1.377 & 4.652 & 1.444 \\
\bottomrule
\end{tabular}
\end{center}

\noindent
As fast as the version written by hand, $3.4$--$4.6\times$ faster than
the formula with \code{inv}; ``in 2 sd'' is the share of test hours whose
value lies within two predictive standard deviations.

\subsection{What PyTorch implements, as formulas}\label{sec:torch}
\file{python/examples/torch_linear.py} (\code{make torch}): a linear layer (\code{torch.nn.Linear}: $Y = XW + \mathbf{1}b\T$) trained
by mini-batch gradient descent on the mean squared error: the forward
pass, the gradients and the optimiser step written as formulas, the whole
loop one ulm.lam function:
\begin{lstlisting}[style=py]
with lam.loop(steps, carry=[W0, b0]) as L:
    W, b = L.carried
    X, Y = batch_x(L.index), batch_y(L.index)
    R = X @ W + ones @ b.T - Y                         # forward, residual
    L.next(W - step * (X.T @ R), b - step * (R.T @ ones))   # SGD step
\end{lstlisting}
Per step ulm.lam generates four kernels: \code{gemm} $XW$ into the buffer
of $-Y$ ($\beta = -1$), the bias as a rank-1 \code{gemm}, \code{gemm}
$X\T R$ with $\alpha = -\mathrm{step}$ and $\beta = 1$ directly into $W$,
and \code{gemv} $R\T\mathbf{1}$ into $b$; $W$ and $b$ stay in place, and
nothing is allocated in the loop.
That needed one more rule: a term $s(AB)$ of a sum with $s$ known at run
time is a product whose \code{gemm} gets $s$ as its $\alpha$ (scales of
scales fold into one, also for run-time scalars, as an \code{arith.mulf}).
Compared with PyTorch~2.15 in float64 on the same BLAS (Accelerate, which
PyTorch uses on macOS; ulm.lam with \code{backend="cblas"}), 100 steps
(seconds):

\begin{center}
\begin{tabular}{rrrrrrr}
\toprule
& & & \multicolumn{2}{c}{ulm.lam} & \multicolumn{2}{c}{PyTorch} \\
\cmidrule(lr){4-5}\cmidrule(l){6-7}
$B$ & $p$ & $k$ & Accelerate & ulmBLAS & autograd & by hand \\
\midrule
32 & 64 & 16 & 0.001 & 0.001 & 0.004 & 0.001 \\
256 & 1024 & 256 & 0.080 & 0.157 & 0.088 & 0.081 \\
1024 & 4096 & 1024 & 3.739 & 9.730 & 3.779 & 3.797 \\
\bottomrule
\end{tabular}
\end{center}

\noindent
``PyTorch autograd'' is \code{nn.Linear} with \code{backward()} and
\code{optim.SGD}, ``by hand'' the same steps as torch operations without
autograd.
ulm.lam is as fast as PyTorch where the \code{gemm}s dominate, and
$5\times$ faster than autograd for a small layer, where PyTorch's work per
operation dominates; the results agree to $10^{-14}$.
So the operations that PyTorch implements in C++ (the forward and backward
of its modules, the optimiser steps) can be written as formulas and
compiled by ulm.lam without losing speed: ulm.lam could implement such a
framework.
Caveats: the gradient formulas are written by hand (ulm.lam does not
differentiate; differentiating its IR symbolically is a natural
extension); the comparison is in float64, while ML mostly uses float32,
which ulm.lam does not support end to end yet; on ulmBLAS, whose
\code{gemm} runs on one thread so far (and without OpenMP here), ulm.lam
is $2$--$2.6\times$ slower; and \code{torch.compile} failed on this
machine (PyTorch's headers do not compile with this Apple clang:
\code{unknown type name 'uint'}).
PyTorch brings its own OpenMP runtime, which aborts the process if
Homebrew's (NumPy's OpenBLAS, ulm.lam's OpenMP) is loaded too; the
showcase therefore runs in a virtual environment with NumPy and PyTorch
from pip (both on Accelerate) and builds ulm.lam without OpenMP.

\subsection{The same in Julia}\label{sec:showcasesjulia}
In Julia (\file{julia/examples/ridge.jl},
\file{attention.jl}, \file{kalman_filter.jl}, \file{bayes.jl}, \code{make julia}), against
what a Julia programmer writes (Julia 1.13 with its OpenBLAS on 4
threads; ulm.lam on ulmBLAS, for attention also on OpenBLAS with 4
threads; seconds, the largest sizes):

\begin{center}
\small
\begin{tabular}{@{}lrrr@{}}
\toprule
& ulm.lam & Julia as written & Julia by hand \\
\midrule
ridge regression, $p = 4000$ & 1.268 & 2.438 & 1.750 \\
Bayesian linear regression, $p = 4000$ & 1.545 & 2.787 & 1.989 \\
Kalman filter, $n = 2000$, 20 steps & 2.655 & 4.811 & 4.465 \\
linear attention, $N = 16384$ (ulmBLAS / OpenBLAS) & 0.0117 / 0.0032 & 0.714 & 0.0021 \\
\bottomrule
\end{tabular}
\end{center}

\noindent
Ridge regression, Bayesian regression and the Kalman filter are
$1.3$--$1.7\times$ faster than the versions written by hand (why Julia's are slower than NumPy's
relative to ulm.lam has not been analysed here).
For attention, ``as written'' is \code{(Q * K') * V}; Julia's \code{*}
with three arguments, \code{Q * K' * V}, already chooses the order of a
chain of three matrices (0.0022\,s), as ulm.lam does.
There ulm.lam on OpenBLAS is as fast as Julia within the spread of these
millisecond timings, and on ulmBLAS it is slower: its \code{gemm} runs on
one thread so far, a hint for ulmBLAS rather than for the mapping.

\noindent
The program needed three general improvements of the compiler.
An \code{assume} keeps the storage of its operand (a symmetric matrix
stored as one triangle stays so, which also removed a completion before
\code{potrf} in ridge regression).
A loop-carried value that is not computed in the buffer of its iteration
argument ($x = Fx$ cannot be in place) is copied into it at the end of the
step, instead of being allocated in every step.
And the reuse of computed products ignores the transposition of
symmetric factors, so that $PH\T$ is found as the transpose of $HP$.
Still possible: $FPF\T$ and $HPH\T$ are symmetric, so \code{gemmtr} would
halve the second product.

\addtocontents{toc}{\protect\setcounter{tocdepth}{1}}
% =====================================================================
\section{Testing strategy}
% =====================================================================

\begin{itemize}
  \item \textbf{Per pass}: \code{lit}/\code{FileCheck} tests on textual IR.
        Properties are part of the types, so tests of inference rules are
        one-liners.
  \item \textbf{Property soundness}: random matrices that satisfy the
        input properties are generated, every inferred property is
        evaluated numerically, and the test fails on any property that
        does not hold.
        This catches wrong inference rules, the most dangerous class of
        bug.
  \item \textbf{Numerical} (implemented): every example is compiled to
        C++ and compared with an independent evaluation.
        \file{examples/cases/} covers the LU solve, plain and transposed,
        upper triangular solves, right solves with $L\inv$ and $L\invT$,
        a Cholesky solve with a vector, a sum of three plain terms and
        $A\inv BC + D$, against naive triple loops and Gaussian
        elimination; all agree to $\le 6\cdot10^{-16}$.
        It also checks that inputs passed by value are unchanged.
  \item \textbf{Structural} (implemented, scripts): kernel selection and
        memory placement are checked on the generated IR, e.g.\ two
        Cholesky factorisations and three copies for GLS, no allocation in
        the loop.
  \item \textbf{Differential}: \code{convert-ulm-lamk-to-linalg} against
        \code{convert-ulm-lamk-to-ulmblas}.
  \item \textbf{Quality}: the LAMP benchmark~\cite{lamp-paper}
        (\url{https://github.com/HPAC/LAMP_benchmark}) and Linnea's 25
        application problems.
        Each derivation is compared against Linnea's FLOP count and
        run time.
\end{itemize}

% =====================================================================
\section{Status and plan}\label{sec:status}
% =====================================================================

\paragraph{Done.}
\begin{itemize}
  \item ODS for both dialects, interfaces and passes, checked by all
        \code{mlir-tblgen} backends, and a kernel catalogue whose TableGen
        syntax is checked;
  \item the \code{ulm_lam} and \code{ulm_lamk} dialects, generated from the
        ODS and loaded into \code{mlir-opt} as a plugin: property
        inference, relaxed type compatibility, verifiers, canonicalisation
        patterns, kernel costs and bufferization;
  \item \code{ulm-lam-normalize}, \code{ulm-lam-infer-properties},
        \code{ulm-lam-select-kernels} with \code{strategy=gmc} including
        inverses, factorisations, solves, $Z\T Z$ and sums
        (\cref{sec:impl}), and \code{ulm-lam-strip-properties};
  \item parameter intents \code{in}, \code{inout}, \code{out}, and a
        generated comment that documents them for the caller;
  \item \code{ulm-lam-results-to-out-params}: matrices by reference, the
        caller owns them ($X=\mathrm{foo}(i)$ becomes
        $\mathrm{foo}(i,X)$, in/out results are updated in place, also by
        external functions via \code{ulm_lamk.call}, inputs passed by value
        are read-only);
  \item \code{ulm-lam-workspace}: temporaries from a caller-provided
        workspace (query, allocate, or check), with queries accumulated over
        calls of generated and external functions;
  \item \code{convert-ulm-lamk-to-ulmblas} for all kernels kernel selection
        emits;
  \item the Python frontend and runtime (\code{from ulm import lam}):
        tracing, loops, external functions as Python callbacks,
        workspaces, NumPy or plain buffers;
  \item the Julia frontend and runtime (\file{julia/UlmLam},
        \cref{sec:julia}): \code{@lam} functions called with Julia arrays,
        \code{@code_lam}, Julia's own BLAS through libblastrampoline,
        callbacks, one workspace per task (calls from several threads);
  \item a second lowering on CBLAS and LAPACK (MKL, OpenBLAS, Accelerate,
        Julia's BLAS);
  \item the operations of the vector spaces $\mathbb{R}^n$ and
        $\mathbb{R}^{m\times n}$, combined with products, transposes,
        solves and explicit inverses, checked against Julia and NumPy;
  \item a Pluto notebook that runs in the browser (Binder, with
        \file{binder/Dockerfile}) and as a static page;
  \item GLS from its normal form, formula (1) in a loop (with a constant
        and with an updated covariance), the loop example and eight kernel
        selection cases, compiled to C++ with
        ulmBLAS by passes alone and verified numerically; the hand-written
        GLS stages 1–4 parse and verify.
\end{itemize}

\paragraph{Compared with Linnea.}
Linnea (\url{https://github.com/HPAC/linnea}, last changed in 2021) was run
on the problems of ulm.lam's examples (\file{examples/linnea}: the inputs in
Linnea's language, \file{run.sh}, and the best algorithm Linnea generates
for each in \file{generated/}; running it needs Python~3.12, TatSu~5.6.1
and setuptools~${<}\,70$, and its \code{--time-limit} option fails).
Linnea finds each solution in under 10\,s.
The kernel sequences (FLOPs for the sizes of the inputs):

\begin{center}
\footnotesize
\begin{tabularx}{\linewidth}{@{}>{\raggedright\arraybackslash}p{0.2\linewidth}>{\raggedright\arraybackslash}X>{\raggedright\arraybackslash}X@{}}
\toprule
problem & Linnea & ulm.lam \\
\midrule
GLS & potrf, trsv, trsm, syrk, potrf, gemv, 2 trsv ($3.55\cdot10^{9}$) & the same \\
GP ($n = 4000$, $m = 1000$) & $\sigma^2$ added to the diagonal, potrf, two trsm ($K\inv K_*$), gemv, gemm for $K_{**} - K_*\T(K\inv K_*)$ ($6.13\cdot10^{10}$) & one trsm ($V = L\inv K_*$), syrk for $K_{**} - V\T V$ ($4.1\cdot10^{10}$) \\
Kalman filter, one step ($n = 1000$, $m = 100$) & $FPF\T$ as potrf($P$), trmm ($FL$), syrk ($2.78\cdot10^{9}$) & two gemm for $FPF\T$ ($4.4\cdot10^{9}$) \\
ridge, Bayesian regression & syrk onto $\lambda I$ (the identity as an array, $\beta = \lambda$), potrf, trsv & the same; the identity is never formed (adddiag) \\
linear attention & $Q(K\T V)$, fixed when the code is generated & the same, chosen at run time \\
SGD step & gemm, \code{ger} for the bias, gemm and gemv with $\alpha = -\mathrm{step}$ & the same, the bias as a gemm with $k = 1$ \\
\bottomrule
\end{tabularx}
\end{center}

\noindent
So in the mapping each has an advantage the other lacks: ulm.lam sees the
Gram structure of the GP covariance (a third fewer FLOPs), Linnea makes
$FPF\T$ symmetric through a Cholesky factorisation of $P$ (without a
factorisation, \code{gemmtr} would bring ulm.lam to $3n^3$), and uses
\code{ger} for a rank-1 product.

\emph{Memory for temporaries} is where the two differ most.
Linnea's code generator (\file{code\_generation/memory/memory.py}) maps
every operand to a ``memory location''; the inputs are the caller's
arrays.
A kernel that overwrites an operand writes into its location if the
operand is not used afterwards (liveness), otherwise the operand is first
copied into a new location (\code{blascopy!}).
Every other result gets a new \code{Array\{Float64\}(undef, ...)},
allocated in every call and freed by Julia's garbage collector.
In the generated code that means:
\begin{center}
\footnotesize
\begin{tabularx}{\linewidth}{@{}>{\raggedright\arraybackslash}p{0.14\linewidth}>{\raggedright\arraybackslash}X>{\raggedright\arraybackslash}X@{}}
\toprule
& Linnea & ulm.lam \\
\midrule
inputs & overwritten as needed: after GLS the caller's $M$, $X$ and $y$ hold factors and solutions, after a Kalman step $F$, $P$, $Q$, $R$, $z$ are destroyed & \code{in}: never written (copied if a kernel must overwrite); \code{inout}: may be overwritten \\
temporaries & \code{Array\{Float64\}(undef, ...)} in every call (GC) & one workspace, with a LAPACK-style size query; given by the caller or allocated once \\
reuse & memory locations with liveness: a result overwrites a dead operand & bufferization's in-place analysis; workspace layout by liveness, branches overlap \\
identity, zero & formed as arrays when a kernel needs them as input (the $C$ of syrk) & never formed (adddiag) \\
results & returned arrays, which may be inputs (the residual $R$ of the SGD step is returned in the buffer of $Y$) & output parameters owned by the caller \\
loops & none (straight-line programs; a filter step called in a loop allocates in every step, and the caller must copy $F$, $Q$, $R$ before each call) & no allocation in a loop; invariant work before it \\
sizes, errors & fixed when the code is generated; Julia exceptions (\code{PosDefException}) & run-time sizes; \code{info} \\
\bottomrule
\end{tabularx}
\end{center}

\noindent
Linnea's choice is reasonable for code that is called once with
throw-away inputs; for a library routine, and for code in loops, it is
what FLENS-LAPACK had to avoid by hand: temporaries inside nested calls,
and arguments that are silently destroyed.

\paragraph{Plan.}
\begin{enumerate}[label=M\arabic*]
  \item \emph{Largely done.} Dialects, verifiers, property closure and
        inference.
        \code{ulm-lam-opt} (\file{tools/ulm-lam-opt.cpp}, mlir-opt with
        the ulm.lam dialects and passes, linked against the plugin) and
        \code{lit}/\code{FileCheck} tests of the passes in \file{test/}
        (\code{make check}, driven by GNU Make, \code{lit} from PyPI).
  \item \emph{Largely done.} \code{ulm-lam-normalize} flattens products
        and sums, pushes transposes to the leaves ($(AB)\T = B\T A\T$,
        also for sums and scales), folds double transposes and inverses
        and scales of scales, and drops identity factors; kernel
        selection inverts products factor by factor and collects the
        scales of a product into the $\alpha$ of its last kernel.
        Open: inverses pushed to the leaves already in the normal form,
        so that other passes see them.
  \item \emph{Done for what kernel selection emits.} \code{ulm_lamk} with DPS,
        bufferization and the ulmBLAS and CBLAS lowerings, \code{trmm}
        included.
  \item \emph{Partly done.} \code{strategy=gmc} with gemm, gemv, solves,
        factorisations and $Z\T Z$.
        \code{syrk} and \code{symm}/\code{symv} with triangles tracked
        (\cref{sec:storage}).
        A triangular factor of a product is a \code{trmm} (the other
        factor copied, then multiplied in place: half the flops of
        \code{gemm}).
        Missing: \code{gemmtr} for symmetric products that are not
        $Z\T Z$,
        the catalogue backend, so that new kernels are added as records,
        and a fallback plan where a factorisation fails
        (\cref{sec:info}).
  \item \code{strategy=search}: sharing across products beyond reuse of
        computed intervals, factorisation choices, anytime search; the LAMP
        benchmark; a full Kalman filter (\cref{sec:kalman}).
  \item \emph{Done:} loop-aware selection: a loop-invariant prefix of a
        product is formed before the loop if the trip count, known at run
        time, makes it pay off (the operator of \cref{sec:kalman}).
        Open: invariant subchains other than a prefix, and a measured
        instead of a fixed weight for matrix--vector flops.
  \item \emph{Done for Python} (\cref{sec:python}): a frontend that
        traces Python to \code{ulm_lam} IR, with a runtime and callbacks.
        Open: \code{strategy=egraph} experiments; QR and $LDL\T$ once
        ulmBLAS has them (\cref{sec:ulmblas}).
\end{enumerate}

\paragraph{Next steps.}
\begin{enumerate}
  \item \emph{Tests}: more \code{lit} tests, so that every rule of
        kernel selection has one (the examples and the frontends' tests
        check the whole pipeline numerically).
  \item \emph{Kernels}: \code{gemmtr} for symmetric products that are
        not $Z\T Z$, \code{ger} for rank-1 products, \code{trmm} for
        triangular factors that are computed (factors of a factorisation)
        and not only parameters.
  \item \emph{Reactant.jl and Enzyme-JAX}: their \code{enzymexla}
        dialect has BLAS and LAPACK operations, and their ``raising''
        recognises linear algebra in lower-level code.
        ulm.lam's kernel selection could take its input from raising and
        emit their operations instead of C++.
  \item \emph{More showcases}: ridge regression, linear attention and the
        Kalman filter are done (\cref{sec:showcases}).
        Bayesian linear regression too, with run-time scalars on products
        and scalar arithmetic.
        Next: \code{gemmtr} for the symmetric products of the Kalman
        filter.
  \item \emph{PyTorch}: PyTorch maps every operation by itself onto
        BLAS/LAPACK or cuBLAS/cuSOLVER, and \code{torch.compile} fuses
        elementwise operations but does no linear algebra mapping
        (\cite{tfpt} measures what is missing).
        Steps: CPU tensors in and out of compiled functions without
        copies (DLPack); a \code{torch.compile} backend that replaces the
        linear algebra in the FX graphs (matrix chains, \code{inv},
        \code{solve}, \code{cholesky}, also in the backward graph of
        autograd) by functions generated by ulm.lam, with the expressions
        of \cite{tfpt} as the showcase; and, separately, float32 and a GPU
        backend (cuBLAS, cuSOLVER), where most PyTorch users compute.
  \item \emph{Mojo} (an idea, to be checked first): Mojo is a
        Python-like language on MLIR whose strength is fast kernels
        (GEMM and others), but which, as far as known here, does no
        linear algebra mapping.
        ulm.lam could serve it as a frontend (formulas written in Mojo,
        the generated library called through Mojo's C interface; whether
        Mojo programs can be traced like Python and Julia ones is open) or
        use Mojo's kernels as a backend, like the CBLAS one (which LAPACK
        functionality Mojo's libraries provide is open).
  \item \emph{Linnea}: the problems of the examples are compared
        (\cref{sec:status}, ``Compared with Linnea'').
        Next from it: \code{gemmtr} and the factorisation trick for
        $FPF\T$, \code{ger} for rank-1 products.
  \item \emph{Linnea's 25 application problems and the LAMP benchmark}
        (\url{https://github.com/HPAC/LAMP_benchmark}): run them through
        ulm.lam and compare FLOPs and coverage per problem with Linnea.
        That gives the precise list of missing rules, factorisations and
        kernels, and the basis for \code{strategy=search}.
\end{enumerate}

\paragraph{Open questions.}
\begin{itemize}
  \item Dynamic sizes: should the compiler use size hints or
        multi-version code with run-time dispatch?
  \item Cost model: FLOPs mispredict memory-bound kernels.
        Measured per-kernel efficiency (\code{cost-model=measured}) is the
        obvious next step.
  \item Should storage formats (packed triangular, banded) become a second
        encoding dimension, or a separate type?
  \item Batched and tensor contractions: should \code{ulm_lam} be
        generalised, or interoperate with \code{linalg}?
\end{itemize}

% =====================================================================
\appendix
\section{Generated code}\label{app:generated}
% =====================================================================

\code{make} compiles every example with \file{scripts/ulm-lam-compile.sh}
and keeps the C++ in \file{examples/*/generated/}.
For every example this appendix shows three stages: the source program
(\code{ulm_lam} IR, written by hand or traced from Python), the kernels
that \code{ulm-lam-select-kernels} chose for it (\code{ulm_lamk} on
tensors, before bufferization: the BLAS and LAPACK calls with their
operands, \file{build/*/*.selected.mlir}), and the generated C++ in full,
as generated (EmitC's variable names, formatted with \code{clang-format}).

\newcommand{\stage}[3]{%
  \par\smallskip\noindent\textbf{#1} (\file{#2}):
  \IfFileExists{\src #2}{%
    \lstinputlisting[style=#3,basicstyle=\ttfamily\scriptsize]{\src #2}}{%
    \fbox{run \code{make} to generate this file}}}
% a program in Python or Julia: the lines from the one containing #3 to the
% one before the one containing #4 (#5: a name for the extracted file)
\newcommand{\appsrc}[6]{%
  \par\smallskip\noindent\textbf{#1} (\file{#2}):
  \IfFileExists{\src #2}{%
    \srcbetween[style=#6,basicstyle=\ttfamily\scriptsize]{\src #2}{#3}{#4}{#5}}{%
    \fbox{missing file}}}
\newcommand{\generated}[5]{%
  \subsection{\texorpdfstring{\code{#1}}{#1}}
  #5
  \stage{Source}{#2}{mlir}
  \stage{Kernels after selection}{#3}{mlir}
  \stage{Generated C++}{#1}{c}}

\generated{examples/gls/generated/gls.cpp}{examples/gls/02-normal.mlir}%
  {build/gls/gls.selected.mlir}{}{Generalised least squares from its
  normal form (\cref{sec:example}).}
\generated{examples/gls/generated/gls_inout.cpp}{examples/gls/gls_inout.mlir}%
  {build/gls/gls_inout.selected.mlir}{}{GLS with $M$ and $y$ declared
  \code{inout} (\cref{sec:ulmblas}).}
\generated{examples/loop/generated/accumulate.cpp}{examples/loop/accumulate.mlir}%
  {build/loop/accumulate.selected.mlir}{}{$X \mathrel{+}= ABC$ in a loop
  (\cref{sec:loops}).}
\subsection{\texorpdfstring{\code{examples/kalman/generated/estimate.cpp}}{examples/kalman/generated/estimate.cpp}}
Formula (1) in a loop with a constant covariance; with enough steps the
operator is formed before the loop (\cref{sec:kalman}).
The same program in Python and in Julia (there with $n = 4000$; the IR
below is written by hand, with $n = 2000$):
\appsrc{Python}{python/examples/kalman.py}{observe = lam.extern}{def best(}{appkalmanpy}{py}
\appsrc{Julia}{julia/examples/kalman.jl}{const observe}{function literal}{appkalmanjl}{py}
\stage{Source}{examples/kalman/estimate.mlir}{mlir}
\stage{Kernels after selection}{build/kalman/estimate.selected.mlir}{mlir}
\stage{Generated C++}{examples/kalman/generated/estimate.cpp}{c}
\generated{examples/kalman/generated/adaptive.cpp}{examples/kalman/adaptive.mlir}%
  {build/kalman/adaptive.selected.mlir}{}{Formula (1) in a loop with the
  covariance updated in every step (\cref{sec:adaptive}).}
\generated{examples/compose/generated/compose.cpp}{examples/compose/compose.mlir}%
  {build/compose/compose.selected.mlir}{}{Functions with workspaces calling
  functions with workspaces (\cref{sec:workspace}).}
\generated{examples/cases/generated/cases.cpp}{examples/cases/cases.mlir}%
  {build/cases/cases.selected.mlir}{}{The kernel selection cases of
  \file{examples/cases/cases.mlir}.}
The showcases are programs in Python and Julia.
For them the appendix shows the Python and the Julia program, the IR the
Python frontend traces from it, the kernels after selection, and the C++
generated from the Python program (the Julia frontend generates the same
kernels on transposed operands, \cref{sec:julia}).

\subsection{\texorpdfstring{\code{python/examples/generated/gp.cpp}}{python/examples/generated/gp.cpp}}
GP regression (\cref{sec:glance,sec:gp}).
\appsrc{Python}{python/examples/gp.py}{@lam.function}{def numpy_literal}{appgppy}{py}
\appsrc{Julia}{julia/examples/gp.jl}{@lam function gp}{function by_hand}{appgpjl}{py}
\stage{IR traced from Python}{build/python/gp_py/gp_py.mlir}{mlir}
\stage{Kernels after selection}{build/python/gp_py/gp_py.selected.mlir}{mlir}
\stage{Generated C++}{python/examples/generated/gp.cpp}{c}

\subsection{\texorpdfstring{\code{python/examples/generated/ridge.cpp}}{python/examples/generated/ridge.cpp}}
Ridge regression (\cref{sec:showcases}).
\appsrc{Python}{python/examples/ridge.py}{@lam.function}{def numpy_literal}{appridgepy}{py}
\appsrc{Julia}{julia/examples/ridge.jl}{@lam function ridge}{literal(X}{appridgejl}{py}
\stage{IR traced from Python}{build/python/ridge_py/ridge_py.mlir}{mlir}
\stage{Kernels after selection}{build/python/ridge_py/ridge_py.selected.mlir}{mlir}
\stage{Generated C++}{python/examples/generated/ridge.cpp}{c}

\subsection{\texorpdfstring{\code{python/examples/generated/attention.cpp}}{python/examples/generated/attention.cpp}}
Linear attention, the order chosen at run time (\cref{sec:showcases}).
\appsrc{Python}{python/examples/attention.py}{@lam.function}{def phi}{appattentionpy}{py}
\appsrc{Julia}{julia/examples/attention.jl}{@lam function attention}{phi(x)}{appattentionjl}{py}
\stage{IR traced from Python}{build/python/attention_py/attention_py.mlir}{mlir}
\stage{Kernels after selection}{build/python/attention_py/attention_py.selected.mlir}{mlir}
\stage{Generated C++}{python/examples/generated/attention.cpp}{c}

\subsection{\texorpdfstring{\code{python/examples/generated/kalman_filter.cpp}}{python/examples/generated/kalman_filter.cpp}}
The Kalman filter, predict and update in a loop (\cref{sec:showcases}).
\appsrc{Python}{python/examples/kalman_filter.py}{measure = lam.extern}{def numpy_literal}{appkalmanfilterpy}{py}
\appsrc{Julia}{julia/examples/kalman_filter.jl}{const measure}{function literal}{appkalmanfilterjl}{py}
\stage{IR traced from Python}{build/python/kalman_filter_py/kalman_filter_py.mlir}{mlir}
\stage{Kernels after selection}{build/python/kalman_filter_py/kalman_filter_py.selected.mlir}{mlir}
\stage{Generated C++}{python/examples/generated/kalman_filter.cpp}{c}

\subsection{\texorpdfstring{\code{python/examples/generated/bayes.cpp}}{python/examples/generated/bayes.cpp}}
Bayesian linear regression, run-time scalars on products (\cref{sec:showcases}).
\appsrc{Python}{python/examples/bayes.py}{@lam.function}{def numpy_literal}{appbayespy}{py}
\appsrc{Julia}{julia/examples/bayes.jl}{@lam function blr}{function literal}{appbayesjl}{py}
\stage{IR traced from Python}{build/python/bayes_py/bayes_py.mlir}{mlir}
\stage{Kernels after selection}{build/python/bayes_py/bayes_py.selected.mlir}{mlir}
\stage{Generated C++}{python/examples/generated/bayes.cpp}{c}

\subsection{\texorpdfstring{\code{python/examples/generated/torch_linear.cpp}}{python/examples/generated/torch_linear.cpp}}
A linear layer trained by gradient descent, compared with PyTorch (\cref{sec:showcases}; no Julia version).
\appsrc{Python}{python/examples/torch_linear.py}{batch_x = lam.extern}{def torch_autograd}{apptorchpy}{py}
\stage{IR traced from Python}{build/python/torch_linear_ulm/torch_linear_ulm.mlir}{mlir}
\stage{Kernels after selection}{build/python/torch_linear_ulm/torch_linear_ulm.selected.mlir}{mlir}
\stage{Generated C++}{python/examples/generated/torch_linear.cpp}{c}

% =====================================================================
\begin{thebibliography}{99}\small
\bibitem{lamp-paper} C.~Psarras, H.~Barthels, P.~Bientinesi.
  The Linear Algebra Mapping Problem. Current state of linear algebra
  languages and libraries. \emph{ACM TOMS}, 2022.
  \href{https://doi.org/10.1145/3549935}{doi:10.1145/3549935}.
\bibitem{linnea-toms} H.~Barthels, C.~Psarras, P.~Bientinesi.
  Linnea: Automatic generation of efficient linear algebra programs.
  \emph{ACM TOMS}, 2021.
\bibitem{linnea-thesis} H.~Barthels.
  \emph{Linnea: A compiler for mapping linear algebra problems onto
  high-performance kernel libraries}. PhD thesis, RWTH Aachen, 2022.
  \href{https://doi.org/10.18154/RWTH-2023-01193}{doi:10.18154/RWTH-2023-01193}.
\bibitem{gmc} H.~Barthels, M.~Copik, P.~Bientinesi.
  The generalized matrix chain algorithm. \emph{CGO}, 2018.
\bibitem{matchpy} M.~Krebber, H.~Barthels.
  MatchPy: Pattern matching in Python. \emph{JOSS}, 2018.
\bibitem{tfpt} A.~Sankaran, N.~A.~Alashti, C.~Psarras, P.~Bientinesi.
  Benchmarking the linear algebra awareness of TensorFlow and PyTorch.
  \href{https://arxiv.org/abs/2202.09888}{arXiv:2202.09888}, 2022.
\bibitem{mom} L.~Chelini, H.~Barthels, P.~Bientinesi, M.~Copik,
  T.~Grosser, D.~G.~Spampinato.
  MOM: Matrix operations in MLIR.
  \href{https://arxiv.org/abs/2208.10391}{arXiv:2208.10391}, 2022.
\bibitem{mlir} C.~Lattner et al.
  MLIR: Scaling compiler infrastructure for domain specific computation.
  \emph{CGO}, 2021.
\bibitem{egg} M.~Willsey et al.
  egg: Fast and extensible equality saturation. \emph{POPL}, 2021.
\bibitem{egglog} Y.~Zhang et al.
  Better together: Unifying Datalog and equality saturation.
  \emph{PLDI}, 2023.
\bibitem{spores} Y.~R.~Wang, S.~Hutchison, J.~Leang, B.~Howe, D.~Suciu.
  SPORES: Sum-product optimization via relational equality saturation
  for large scale linear algebra. \emph{PVLDB}, 2020.
\bibitem{eqsat} eqsat: An equality saturation dialect for non-destructive
  rewriting. \emph{EGRAPHS}, 2025.
  \href{https://arxiv.org/abs/2505.09363}{arXiv:2505.09363}.
\bibitem{dialegg} DialEgg: Dialect-agnostic MLIR optimizer using equality
  saturation with egglog. \emph{CGO}, 2025.
  \href{https://doi.org/10.1145/3696443.3708957}{doi:10.1145/3696443.3708957}.
\bibitem{flens} M.~Lehn, A.~Stippler, K.~Urban.
  FLENS -- A flexible library for efficient numerical solutions.
  \emph{Proceedings of Equadiff 11}, 2005.
  \url{https://www.dml.cz/bitstream/handle/10338.dmlcz/700445/Equadiff_11-2005-2_55.pdf}.
\bibitem{flens-github} M.~Lehn.
  FLENS: Flexible Library for Efficient Numerical Solutions.
  \url{https://github.com/michael-lehn/FLENS}.
\bibitem{ulmblas} M.~Lehn.
  ulmBLAS: a BLAS implementation in C++.
  \url{https://github.com/michael-lehn/ulmBLAS}.
\bibitem{ulmblasjl} M.~Lehn.
  UlmBLAS.jl: Julia interface to ulmBLAS.
  \url{https://github.com/michael-lehn/UlmBLAS.jl}.

\bibitem{lapis} LAPIS: A performance portable, high productivity compiler
  framework. 2025.
  \href{https://arxiv.org/abs/2509.25605}{arXiv:2509.25605}.
\end{thebibliography}

\end{document}
